\documentclass[12pt,letterpaper]{article}

% ---- Packages ----
\usepackage[utf8]{inputenc}
\usepackage[T1]{fontenc}
\usepackage{amsmath,amssymb}
\usepackage{newtxtext,newtxmath}  % Times New Roman text + math (replaces mathptmx)
\usepackage[margin=1in]{geometry}
\usepackage{setspace}
\usepackage{titlesec}
\usepackage{enumitem}
\usepackage{xr-hyper}
\usepackage{hyperref}
\externaldocument{1 - General Overview}
\externaldocument{3 - Gravity}
\externaldocument{4 - Magnetic and Electric Fields}
\externaldocument{5 - Strong and Weak Forces}
\usepackage{microtype}
\usepackage{booktabs}

% ---- Section formatting ----
\titleformat{\section}{\large\bfseries\sffamily}{\thesection.}{0.5em}{}
\titleformat{\subsection}{\normalsize\bfseries\sffamily}{\thesubsection}{0.5em}{}
\titleformat{\subsubsection}{\normalsize\bfseries}{\thesubsubsection}{0.5em}{}

% ---- Spacing ----
\setstretch{1.15}
\setlength{\parskip}{0.5em}
\setlength{\parindent}{0em}

% ---- Custom commands ----
\newcommand{\rhoa}{\rho_{\!A}}
\newcommand{\Ka}{K_{\!A}}
\newcommand{\vin}{v_{\mathrm{in}}}
\newcommand{\rs}{r_s}
% \hbar is already defined in LaTeX; no redefinition needed

\begin{document}

\begin{center}
{\LARGE\bfseries Unifying Theory of Dynamic Aether Mechanics:\\
Quantum Mechanics}\\[1em]
{\large Erik Otto}
\end{center}

\vspace{1em}

\begin{abstract}
Within the Dynamic Aether Mechanics framework, quantum phenomena emerge from
the properties of a quantum superfluid with a logarithmic equation of state.  This
paper derives wave--particle duality from extended condensate excitations (photons
as spin-wave modes of the spinor condensate, with Derrick's theorem forbidding
stable localized particles in free space) that condense near matter via the Dynamic
Casimir Effect.  Energy quantization arises from angular phase-matching conditions
in the atomic resonant cavity formed by the nuclear flow gradient, reproducing
the hydrogen energy levels $E_n = -13.6/n^2$~eV and all spectral series.  The Pauli
exclusion principle follows from topological vortex interactions: same-circulation
vortices cannot share a spatial mode because their superposed velocity fields
carry divergent energy.  The Schr\"odinger equation is recovered via the Madelung
inverse construction with Wallstrom's topological quantization condition resolved
by the quantized circulation of the medium.  The electron is identified as a
quantized vortex ring with mass determined by $m_e c^2 = \tfrac{1}{2}\rho_A
\kappa_e^2 r_{\mathrm{loop}} [\ln(8r_{\mathrm{loop}}/\xi) - \tfrac{1}{2}]$.
\end{abstract}

\vspace{1em}

\bigskip
\noindent\textit{This paper is part of a series presenting the Unifying Theory of
Dynamic Aether Mechanics.  For the foundational concepts of Aether binding,
the three independent properties (density, pressure, temperature), and the
unification of force constants, see Paper~1: General Overview \cite{otto-overview}.}
\bigskip

\section{Wave--Particle Duality}
% ============================================================

\subsection{Photons as Spin-Wave Excitations of the Condensate}

The quantum superfluid framework developed in Section~\ref{sec:qss} establishes
the Aether as a spinor condensate with order parameter
$\Psi = (\psi_\uparrow, \psi_\downarrow)$.\ This spinor structure supports
spin-wave excitations: modes in which the two components oscillate in
antiphase ($\delta\psi_\uparrow = -\delta\psi_\downarrow$), modulating the
local spin orientation of the condensate without changing the total density.\
Photons are identified with these spin-wave modes.\ At macroscopic wavelengths
(all observable light), spin-orbit coupling produces effective mass currents
from the spin oscillation, making the photon indistinguishable from a
first-sound perturbation; the effective propagation speed is therefore the
two-fluid first-sound speed $c_1$ (Paper~6~\cite{otto-cosmology},
Section~5.1), which is relevant for dark matter lensing consistency.

Two transverse polarization states arise because the spin-wave mode carries a
vector degree of freedom: the direction of the spinor oscillation in the internal
$(\psi_\uparrow, \psi_\downarrow)$ space.\ This direction projects onto two
independent components, corresponding to the two physical polarizations of light.\
In laboratory spinor BECs, these spin-wave excitations are well-studied and
propagate as gapless Goldstone modes when the system has continuous spin-rotation
symmetry \cite{stamper-kurn2013}.

The spatial transversality of these modes---the requirement that the oscillation
direction is perpendicular to the propagation direction---follows from spin-orbit
coupling in the condensate.\ In superfluid helium-3 (the A-phase), Volovik showed
that spin-orbit locking ties the internal spinor direction to the spatial orientation
of the order parameter, producing emergent transverse gauge bosons from the condensate's
Goldstone modes \cite{volovik2003}.\ The same mechanism applies here: the Aether
condensate's spinor structure, combined with the spin-orbit interaction that relates
the internal spin direction to the spatial wave vector, ensures that the spin-wave
excitations are transverse.

The spin-wave dispersion relation is $\omega = c\,k$---exactly non-dispersive: all
frequencies travel at the same speed $c$.\ This is confirmed to extraordinary
precision by the Fermi LAT gamma-ray timing constraints, which constrain
Lorentz-violating photon dispersion to energy scales above
$7.6\,E_{\mathrm{Planck}}$ \cite{vasileiou2013}.\ Since photons propagate in the
spin-wave channel, this constrains the \emph{spin} healing length $\xi_s \ll \xi$,
consistent with the spin-wave channel being vastly stiffer than the density channel
(Section~\ref{sec:strong-constraints}).

The contrast with gluons is instructive.\ Photons propagate through \emph{unbound}
Aether, whose response is confirmed to be perfectly linear (Hookean) to $10^{-8}$
precision \cite{reinhardt2007}.\ Gluons are excitations of \emph{compressed} bridge
material---a dense condensate strand whose nonlinear properties produce the
self-interactions that distinguish the strong force from electromagnetism.\ The
distinction is one of medium: the linear bulk condensate versus the nonlinear
compressed strand.

\textbf{Connection to the flow-vector description.}\quad
Paper~4~\cite{otto-em} describes the electric field as the Aether ``flow vector''
and the magnetic field as its vorticity---a macroscopic velocity-field picture.\
This is not in conflict with the spin-wave identification.\ In superfluid
helium-3, Volovik \cite{volovik2003} showed that spin-wave Goldstone modes,
through spin-orbit coupling, produce effective mass currents indistinguishable
from velocity-field perturbations at macroscopic scales.\ The same mechanism
applies here: the internal spinor oscillation (the microscopic reality) couples
to the spatial velocity field through spin-orbit locking, producing the flow
patterns that Paper~4 describes as the electric and magnetic fields.\ The
``flow vector'' is the macroscopic effective description of what is
microscopically a spin-wave excitation.

A quantitative derivation showing that the spin-wave dispersion relation reproduces
Maxwell's equations in full---including the transversality condition
$\mathbf{k} \cdot \mathbf{E} = 0$ and the absence of a longitudinal photon---requires
specifying the spin-orbit coupling Hamiltonian of the Aether condensate and remains
an open calculation.


\subsection{Why Photons Cannot Be Localized Particles in Free Space}

The Derrick scaling analysis (Section~\ref{sec:derrick}) proves that no stable, localized
three-dimensional field configuration exists in the logarithmic condensate.\ Under uniform
rescaling $\psi(\mathbf{r}) \to \psi(\mathbf{r}/L)$, the energy of any localized
perturbation scales as $E(L) = AL + BL^3$ with $A, B > 0$.\ Since this is monotonically
increasing, the configuration always lowers its energy by shrinking to zero size---there
is no stable soliton.

This result is a \emph{feature}, not a limitation.\ It explains why photons propagate as
extended waves rather than as compact particles: the unbound Aether simply does not support
stable localized excitations.\ A ``photon particle'' in free space would be an unstable
configuration that immediately disperses back into a wave.\ Stable localized structures
(electrons, quarks) require a different topological mechanism---quantized circulation---which
photons lack because they carry no topological charge.

\subsection{Energy Quantization: $E = h\nu$}

The condensate obeys bosonic commutation relations
$[\hat{a}, \hat{a}^\dagger] = 1$ (Section~\ref{sec:qss}), permitting unlimited occupation
of each mode but requiring that each mode's energy comes in discrete quanta.\ The quantum
of action $h$ that sets the vortex circulation quantum $\kappa_0 = h/m_A$ also sets the
minimum energy per wave mode: $E = h\nu$.\ Energy quantization of light is not an
independent postulate---it follows from the same quantum structure of the medium that
produces circulation quantization and the Pauli exclusion principle.

\subsection{Wave-to-Particle Transition Near Matter}
\label{sec:wave-to-particle}

If photons cannot exist as stable particles in free space, what happens when light is
detected?\ Every atom generates a radial Aether inflow (its electric field;
Section~\ref{sec:electric-field}), creating a velocity gradient that is strongest near the
nucleus.\ When a propagating wave enters this strong-inflow environment:
%
\begin{enumerate}[leftmargin=2em]
  \item The inflow gradient refracts the wave inward, compressing it spatially---analogous
    to how sound is bent and concentrated by wind shear in the atmosphere.
  \item The local energy density increases as the same wave energy is concentrated into a
    smaller volume.
  \item The rapidly moving vortex boundaries of atomic electrons provide the conditions for
    the Dynamic Casimir Effect: fluctuations in the wave are amplified at the oscillating
    boundary, facilitating the conversion of distributed wave energy into a localized
    excitation.
\end{enumerate}
%
Detection is therefore a physical process: the wave condenses into a localized energy
deposit at a specific atom.\ This is not a ``wavefunction collapse'' requiring an external
postulate---it is the natural behavior of a spin-wave excitation entering a strong-inflow region
bounded by oscillating vortex structures.\ (The Derrick analysis of
Section~\ref{sec:derrick} applies to the \emph{uniform} condensate; the atomic
environment's steep flow gradient and oscillating vortex boundaries modify the
energy functional, enabling localized energy transfer that is forbidden in free space.)

\subsection{The Double-Slit Experiment and Its Variants}

The double-slit experiment is resolved without invoking complementarity or
observer-dependent collapse.\ Between source and screen, the photon is a
spin-wave excitation obeying exact superposition---it passes through both
slits and interferes with itself, identically to any classical wave in a
linear medium.\ At the screen, each atom's strong-inflow environment condenses
the wave energy into a localized detection event (probability $\propto$
$|\psi|^2$), building up the interference pattern one event at a time.\
Which-path detection destroys the pattern not because ``observation collapses
the wave'' but because the detector's vortex-boundary interactions extract
energy and disrupt phase coherence.

This framework naturally accommodates all major variants: single-photon
interference \cite{taylor1909,grangier1986} (each excitation interferes with
itself); massive-particle interference up to 25{,}000~amu molecules
\cite{fein2019} (all particles are condensate excitations with de~Broglie
wavelength $\lambda = h/(mv)$); Wheeler's delayed-choice experiment
\cite{wheeler1978,jacques2007} (the wave always propagates through both
arms---no retrocausality, just delayed boundary conditions);
interaction-free measurement \cite{elitzurvaidman1993,kwiat1995} (absence
of a wave component eliminates destructive interference); and the continuous
$V^2 + D^2 \leq 1$ complementarity tradeoff \cite{englert1996} (partial
scattering proportionally disrupts phase coherence).\ Decoherence of large
molecules follows from thermal photon emission carrying which-path
information, with no fundamental cutoff.


\subsection{Entanglement-Dependent Experiments}

Two further double-slit variants---the quantum eraser and the delayed-choice
quantum eraser---involve entangled particle pairs.\ Their explanation within the
Aether framework draws on the non-separability of correlated condensate
excitations (Section~\ref{sec:entanglement-binding}): because both members of an
entangled pair are topologically bound features of a single quantum condensate,
their joint state cannot be factored into independent descriptions.\ The
experimental results follow naturally from this wave-based picture.

\textbf{The quantum eraser} (Scully and Dr\"uhl 1982 \cite{scullydruhl1982}).\ Photons
passing through a double slit are tagged with which-path information by coupling
each path to an orthogonal polarization state (or an auxiliary quantum system).\
With the tags intact, no interference is observed---orthogonal Aether wave modes
do not superpose.\ However, a subsequent measurement that projects the tagging
system onto a basis that mixes the two tags (``erasing'' the which-path
information) restores interference in the corresponding sub-ensemble.\ In the
Aether model, the tagged waves are spin-wave modes with orthogonal
polarization structures (distinct patterns of the spinor order parameter).\ Orthogonal modes do not interfere, just as orthogonally polarized classical
waves produce no fringe pattern.\ The eraser measurement reprojects the wave
onto a shared polarization basis, recovering the coherence needed for
interference in that subset of events.\ The total pattern (summed over all
tagging outcomes) never shows fringes, because the fringe and anti-fringe
sub-ensembles cancel exactly.

\textbf{The delayed-choice quantum eraser} (Kim et~al.\ 2000
\cite{kim2000}).\ Entangled photon pairs are produced via spontaneous parametric
down-conversion.\ The ``signal'' photon passes through a double slit and is
detected at a screen; the ``idler'' photon travels a longer path to one of
several detectors.\ Depending on which detector registers the idler, the
which-path information is either preserved or erased.\ The idler detection
can occur \emph{after} the signal photon has already been recorded.

The total signal-photon pattern shows no interference---it is featureless.\
But when the data is sorted by idler outcome, the sub-ensemble whose idlers
reached the erasure detectors reveals interference fringes (and anti-fringes),
while the sub-ensemble whose idlers reached the which-path detectors shows
none.

In the Aether framework, the entangled pair originates as two correlated
spin-wave excitations produced simultaneously in the down-conversion process.\
The correlations---encoded in the joint polarization and phase structure of
the two waves---are established at the source, not imposed retroactively by
measurement.\ The signal wave at the screen carries these correlations as
substructure that is invisible in the total intensity pattern but becomes
apparent when the data is sorted by the idler's outcome.\ The idler
measurement does not \emph{change} what happened to the signal; it
\emph{sorts} the signal events into pre-existing correlated sub-ensembles.\
No faster-than-light influence is required, and no retrocausality is invoked.

\subsubsection{Entanglement as Condensate Non-Separability}

\label{sec:entanglement-binding}
The preceding explanations of the quantum eraser and delayed-choice quantum eraser
rest on the claim that correlated wave pairs maintain a joint quantum state that
cannot be factored into independent single-wave descriptions.\ Within the Aether
framework, this non-separability is not an additional postulate---it follows
directly from the binding principle established in Section~\ref{sec:core-binding}.

The core axiom of the theory is that all matter and radiation are excitations of a
single quantum superfluid condensate described by a spinor order parameter
$\Psi = (\psi_\uparrow, \psi_\downarrow)$.\ Particles are not objects placed
\emph{in} the Aether; they are topological deformations \emph{of} it---vortex
rings, compression solitons, and bridge structures that cannot be separated from
the condensate.\ This binding is what gives them mass, confinement, and stability.

The same binding principle that explains confinement also addresses entanglement.\
When two correlated excitations are produced---for example, a photon pair from
spontaneous parametric down-conversion, or a particle--antiparticle pair from
vacuum fluctuations---they are not two independent objects that happen to share
information.\ They are two correlated features of a single quantum field that were
never separate to begin with.\ The condensate's wavefunction is one state, not two,
and the correlations between excitations are encoded in that global state from the
moment of creation.

\textbf{Why this resolves the Bell inequality challenge.}\ Bell's theorem
\cite{bell1964} proves that no theory based on \emph{local hidden variables} can
reproduce the correlations predicted by quantum mechanics.\ The key assumption in
Bell's proof is \emph{separability}: that the joint probability of outcomes at two
detectors factors into independent probabilities, each depending only on its local
hidden variable and the local detector setting.\ A classical fluid with
pre-determined wave properties would satisfy this assumption and would therefore
obey Bell's inequality.

The Aether condensate is not a classical fluid.\ It is a quantum superfluid whose
excitations obey bosonic commutation relations $[\hat{a}, \hat{a}^\dagger] = 1$
(Section~\ref{sec:qss}).\ Two correlated excitations of this condensate share a
joint quantum state that does not factor into independent local descriptions---the
separability assumption fails.\ This is not a signal propagating between the
excitations, nor a hidden variable carried by each; it is the irreducible quantum
coherence of the medium itself.

\textbf{BEC pair production as a direct analogue.}\ This mechanism is not
speculative---it is observed in laboratory Bose--Einstein condensates.\
Bogoliubov excitations of a BEC naturally arise in correlated momentum pairs
$(k, -k)$ that are entangled by construction \cite{steinhauer2016}.\ These pairs
violate classical correlation bounds precisely because they are excitations of
a single quantum field, not independent classical objects.\ In the Aether
framework, photon pair production (SPDC) is the cosmological-scale analogue of
Bogoliubov pair production: the same quantum statistics that produce entangled
phonon pairs in a laboratory BEC produce entangled photon pairs in the Aether
condensate.

\textbf{No-signaling is automatic.}\ Because the total detection pattern at either
detector is completely featureless without post-selection---the marginal
distribution at each detector is uniform regardless of what happens at the
other---no information is transmitted faster than light.\ The non-local
correlations are only revealed when results from both detectors are compared after
the fact.\ This is consistent with the theory's requirement that all causal
influences propagate at or below the local speed of sound $c_s = \sqrt{\Ka/\rhoa}$ (which equals $c$ at ambient density).

The Aether framework thus provides a physical picture of entanglement: it is the
quantum coherence of the condensate ground state, inherited by excitation pairs
through the binding that ties all matter and radiation to the same underlying
field.\ The same mechanism that confines quarks inside protons (topological binding
to the condensate) also explains why entangled photons violate Bell inequalities
(they are correlated excitations of the quantum field they are bound to).\ A formal
derivation showing that the condensate's field operators produce exactly the
Bell-violating correlations predicted by quantum mechanics remains an important
goal for future work.


\subsection{The Photoelectric Effect}

The photoelectric effect ($E = h\nu$, threshold at the work function) is
retained with one change in physical picture: the incoming spin-wave
excitation condenses into a localized quantum at a single surface atom
through the DCE mechanism described in Section~\ref{sec:wave-to-particle}.\ The instantaneous
response follows because the wave energy is already distributed across
the surface---condensation occurs on atomic timescales, not classical
accumulation timescales.



% ============================================================
\section{Energy Quantization and Atomic Spectra}
% ============================================================

The preceding sections establish that the Aether is a physical medium with density
$\rhoa$, bulk modulus $\Ka$, and a wave speed $c = \sqrt{\Ka/\rhoa}$ that is constant
in the free medium.\ The electric field is a radial Aether inflow, and every atom
generates a structured velocity environment through this inflow.\ A natural question
follows: if the Aether is a real medium that supports waves, and if atoms create
structured flow environments in this medium, then atoms should function as resonant
cavities---and their energy levels should be discrete for the same reason that a
vibrating drumhead or an organ pipe supports only certain frequencies.

This section develops that connection.

\subsection{The Atomic Resonant Cavity}

Every nucleus generates a radial Aether inflow (its electric field), creating a flow
velocity gradient that is strongest near the nucleus and weakens with distance.\ This
velocity gradient constitutes a refractive environment: waves propagating tangentially
are refracted inward by the radial flow shear, analogous to how sound is bent by wind
gradients in the atmosphere.

An electron in this framework is a permanent Aether binding---a spinning vortex that is
itself a wave-like disturbance in the medium.\ When an electron is bound to a nucleus,
its wave structure must be self-consistent within the density environment created by the
nucleus.\ Specifically, the electron's wave must satisfy two conditions:

\begin{enumerate}[leftmargin=2em]
  \item \textbf{Angular phase matching.}\ After propagating around one complete orbit, the
    wave must return to the same phase---it must constructively interfere with itself.\
    If it does not, destructive interference eliminates the mode.\ This requires an integer
    number of wavelengths around the orbital path:
    %
    \begin{equation}
      n\lambda = 2\pi r
    \end{equation}
    %
    where $n$ is a positive integer and $\lambda$ is the de~Broglie wavelength of the
    electron in the Aether medium.\ This is identical to the condition that selects
    resonant modes on a vibrating ring or in a whispering gallery: the wave must
    close on itself.

  \item \textbf{Radial confinement.}\ The inflow gradient creates an effective boundary.\
    Near the nucleus, the strong radial flow refracts the electron wave inward---confined.\
    At large distances, the inflow weakens and the confining effect vanishes.\ Only certain
    radial profiles produce standing waves that remain bound (decaying at large~$r$) rather
    than radiating away into the medium.\ Non-resonant radial modes leak energy into the
    surrounding Aether and cannot persist.
\end{enumerate}

Together, these two conditions---angular phase matching and radial confinement---select
a discrete set of resonant modes from the continuous spectrum of possibilities.\ Each mode
is labeled by integers: the principal quantum number~$n$ (total number of nodes), the
angular quantum number~$l$ (angular nodes), and the magnetic quantum number~$m$ (orientation
of the angular pattern).\ This is directly analogous to the modes of a vibrating drumhead,
which are labeled by the number of radial nodes and the number of angular nodal lines, and
whose allowed frequencies are determined by the zeros of Bessel functions.

The critical distinction from standard quantum mechanics is that in the Aether model, the
boundary conditions have a physical origin: the flow structure of the medium creates the
cavity.\ The electron wave is not an abstract probability amplitude---it is a real
disturbance in a real medium, confined by the refractive properties of that medium.

\subsection{Derivation of Energy Levels}

The energy levels of hydrogen can be derived from the resonance condition combined with the
Coulomb force balance.\ The electron, treated as a wave in the Aether medium, has a
de~Broglie wavelength determined by its momentum:
%
\begin{equation}
  \lambda = \frac{h}{m_e v}
\end{equation}
%
The angular phase-matching condition $n\lambda = 2\pi r$ gives the
quantization of angular momentum, $m_e v r = n\hbar$---not as an axiom,
but as a consequence of requiring constructive self-interference of a wave
propagating in the Aether around a closed path.\ Combining this with the
Coulomb force balance $ke^2/r^2 = m_e v^2/r$ yields quantized radii
$r_n = a_0 n^2$ (where $a_0 = \hbar^2/(m_e k e^2) = 0.529$~\AA{} is the
Bohr radius), and orbital velocities $v_n = \alpha c / n$, where
$\alpha = ke^2/(\hbar c) \approx 1/137$.\ The fine structure constant thus
acquires a direct physical meaning: it is the ratio of the electron's
orbital velocity to the wave speed of the medium, analogous to a Mach
number.\ That $\alpha \ll 1$ confirms that atomic physics is deeply
non-relativistic---the electron's motion is a small perturbation of the
Aether.

The total energy of each resonant mode is:
%
\begin{equation}
  E_n = -\frac{1}{2}\,\alpha^2\,m_e c^2 \cdot \frac{1}{n^2}
  = -\frac{13.6\;\mathrm{eV}}{n^2}
\end{equation}
%
This result connects three quantities that already have Aether interpretations in this
theory: $c = \sqrt{\Ka/\rhoa}$ (wave speed), $\alpha$ (electron flow strength relative to
the medium), and $m_e c^2$ (the rest energy of the permanent binding).\ The ground state
energy of hydrogen is not an independent parameter---it is determined by the same medium
properties that govern electromagnetism and wave propagation.

\subsection{Consistency with Standard Atomic Physics}

The resonant cavity model reproduces all standard results of atomic
spectroscopy without modification.\ The virial theorem ($T = -E$,
$V = 2E$) follows from the $1/r$ potential.\ Spectral lines match the
Rydberg formula $R_\infty = \alpha^2 m_e c/(2h) = 1.0974 \times
10^7$~m$^{-1}$, giving Lyman-$\alpha$ at 121.6~nm, H-$\alpha$ at
656.3~nm, and the Lyman limit at 91.2~nm ($= 13.6$~eV).

The cavity model reproduces all standard results of atomic spectroscopy
without modification: selection rules (from angular momentum exchange
between the medium and the transitioning mode), fine structure (from
spin-orbit coupling of the electron vortex with its own Aether flow),
hyperfine structure including the 21-cm line at 1420.405~MHz (from
coupling between electron and proton vortex flows), Zeeman and Stark
effects (from external Aether flow perturbations splitting the cavity
modes), multi-electron shell filling, the quantum defect formula, and
exotic systems (positronium, muonic hydrogen).\ All follow from the
Aether resonance conditions without additional assumptions.

% ============================================================
\section{The Pauli Exclusion Principle: Topological Vortex Exclusion}
\label{sec:pauli}
% ============================================================

One of the deepest results in quantum mechanics is the Pauli exclusion principle:
no two identical fermions (half-integer spin particles) can simultaneously occupy
the same quantum state.\ For electrons in atoms, this means no two electrons can
share all four quantum numbers ($n$, $l$, $m_l$, $m_s$).\ The exclusion principle
governs the structure of the periodic table, the stability of matter, and the
properties of every material substance.

Despite its foundational importance, the spin-statistics theorem---which connects
half-integer spin to antisymmetric wavefunctions---has resisted elementary
explanation.\ Feynman remarked in 1963 that he could not provide one.\ The proofs
of Fierz (1939) and Pauli (1940) require abstract axioms (Lorentz invariance,
positive energy, microcausality), and Duck and Sudarshan's comprehensive review
\cite{duck1997} concluded that the connection between spin and statistics remains
one of the least intuitively understood results in physics.

The Aether vortex model offers a physical picture: the exclusion principle arises
from the \emph{topological instability of multiply quantized vortices}.

\subsection{Electrons as Topological Vortices}

In this theory, an electron is a persistent spin vortex in the Aether---a
circulatory flow pattern with quantized angular momentum $\pm\hbar/2$.\ The
vortex has a well-defined \emph{phase winding}: the Aether flow completes a $2\pi$
phase rotation around the vortex core.\ This $2\pi$ winding is the defining
topological characteristic of a spin-$\tfrac{1}{2}$ object.\ In the two-fluid
model (Paper~6~\cite{otto-cosmology}), the electron vortex is a quantized
circulation in the \emph{superfluid condensate}; all particle structure, binding,
and electromagnetic phenomena are condensate dynamics.\ The normal component
(thermal excitations) does not support stable vortices and is relevant only for
the cosmological thermal structure discussed in Paper~6.

Two fundamental properties of vortex topology are relevant:

\begin{enumerate}[leftmargin=2em]
  \item \textbf{Topological charge conservation.}\ In any medium supporting
    quantized vortices, the total topological charge (net winding number) is
    conserved.\ Vortices are created and destroyed only in pairs of opposite
    winding.\ This is confirmed experimentally in superfluid helium and
    Bose--Einstein condensates \cite{matthews1999}.

  \item \textbf{The double-cover property.}\ A spin-$\tfrac{1}{2}$ vortex
    requires a $4\pi$ (720$^\circ$) rotation to return to its original state; a
    $2\pi$ rotation produces a sign change.\ This is the physical content of the
    SU(2) double cover of SO(3), demonstrated by Dirac's belt trick: a vortex
    tethered to the surrounding medium by its flow field cannot be untwisted by a
    single $360^\circ$ rotation, but can be untwisted by $720^\circ$.\ In the
    Aether model, electron vortices are literally tethered to the medium through
    their circulation, giving the belt trick a direct physical realization.
\end{enumerate}

\subsection{Why Same-Spin Vortices Cannot Share an Orbital}

Consider two electron vortices placed in the same resonant cavity mode (same $n$,
$l$, $m_l$) with the \emph{same} spin orientation (same $m_s$---same circulation
direction).\ Their phase windings superpose:
%
\begin{equation}
  \text{Same spin: } 2\pi + 2\pi = 4\pi \quad\text{(net winding)}
\end{equation}
%
A $4\pi$ phase winding is topologically equivalent to a spin-1 object---a boson,
not two fermions.\ The system has undergone a topological transition that changes
its statistical character.\ This is physically forbidden for two reasons:

\textbf{Energetic instability.}\ A vortex with winding number $s$ has kinetic
energy proportional to $s^2$, while $s$ individual unit vortices have total energy
proportional to $s$.\ Therefore a doubly quantized vortex ($s = 2$, energy
$\propto 4$) is energetically unstable relative to two singly quantized vortices
(total energy $\propto 2$).\ This instability is well established in superfluid
physics: multiply quantized vortices in Bose--Einstein condensates spontaneously
decay into singly quantized vortices \cite{shin2004}.

\textbf{Topological inconsistency.}\ Two spin-$\tfrac{1}{2}$ excitations cannot
smoothly deform into a single spin-1 excitation while conserving both angular
momentum and topological charge.\ The transition would require either the creation
or destruction of a net $2\pi$ winding---violating topological charge conservation.

\subsection{Why Opposite-Spin Vortices Can Share an Orbital}

Two electron vortices with \emph{opposite} spin orientations (opposite circulation
directions) have phase windings that partially cancel:
%
\begin{equation}
  \text{Opposite spin: } +2\pi + (-2\pi) = 0 \quad\text{(net winding at overlap)}
\end{equation}
%
The combined configuration retains zero net topological charge at the overlap
region.\ Each vortex maintains its individual identity and topological character---the
system remains two distinct spin-$\tfrac{1}{2}$ objects.\ This is the only stable
two-electron configuration within a single resonant cavity mode, and it explains
the fundamental rule of atomic physics: each orbital holds at most two electrons,
which must have opposite spins.

\subsection{The Exchange Phase from Vortex Topology}

Berry and Robbins (1997) \cite{berry1997} demonstrated that the exchange phase
factor $(-1)^{2S}$ for identical particles arises as a \emph{geometric
(Berry) phase} of topological origin, associated with non-contractible paths in
the configuration space of identical particles.

In the Aether model, this result has a direct physical interpretation.\ When two
identical electron vortices are exchanged by moving them around each other, the
Aether flow between them acquires a twist.\ Because each vortex is a
spin-$\tfrac{1}{2}$ object requiring $4\pi$ rotation to return to its original
state, a single exchange (equivalent to a $2\pi$ relative rotation) maps the
two-vortex state to its negative:
%
\begin{equation}
  \Psi(1,2) \;\xrightarrow{\text{exchange}}\; -\Psi(2,1)
\end{equation}
%
This antisymmetry under exchange is not imposed as an axiom but \emph{follows
from the topology of vortex flow fields in the three-dimensional Aether}.\ If
the two vortices are in the same quantum state, $\Psi(1,2) = \Psi(2,1)$, and
the antisymmetry condition gives $\Psi = -\Psi = 0$---the state cannot exist.

The dependence of exchange statistics on spatial dimensionality provides further
support: in two dimensions, the configuration space of identical particles has the
braid group $B_N$ as its fundamental group rather than the symmetric group $S_N$,
allowing continuous exchange phases (anyons).\ Leinaas and Myrheim (1977)
\cite{leinaas1977} predicted this, and anyonic statistics have been experimentally
confirmed in fractional quantum Hall systems.\ In three dimensions, only $\pm 1$
exchange phases are possible---exactly the boson/fermion dichotomy.\ The fact that
statistics depends on dimensionality is a strong indication that its origin is
\emph{topological}, not dynamical.

\subsection{Consequences of Vortex Exclusion}

The topological exclusion principle reproduces the full structure of atomic
and condensed-matter physics.\ Shell capacities ($2n^2$) follow from the
two-vortex-per-mode limit summed over $l = 0, \ldots, n{-}1$, and Hund's
rules follow from minimizing vortex flow-field overlap.\ Electron
degeneracy pressure ($P \propto n_e^{5/3}$) arises because filling all
low-energy modes forces additional electrons into higher modes; the
Chandrasekhar limit ($\sim\!1.44\,M_\odot$ \cite{chandrasekhar1931}) marks
where vortex flow speeds become relativistic.\ Covalent bonding is the
sharing of vortex modes between adjacent atomic cavities, with
bonding/antibonding configurations from constructive/destructive
flow-field overlap.

The vortex mechanism is directly confirmed in superfluid experiments:
bosonic $^4$He condenses without pairing (no exclusion), while fermionic
$^3$He requires Cooper pairing at 2.491~mK \cite{osheroff1972}; multiply
quantized vortices spontaneously split into singly quantized ones
(confirming $s^2$ instability) \cite{shin2004}; and half-quantum vortices
in $^3$He \cite{autti2016} directly realize the spin-$\tfrac{1}{2}$
topology.



% ============================================================
\section{Constraints from Quantum Superfluid Structure}
\label{sec:qss}
% ============================================================

The topological vortex model of the Pauli exclusion principle carries a profound
implication: the Aether is not merely a classical elastic medium but a
\emph{quantum superfluid} with specific mathematical properties.\ This section
traces the constraints that spin-statistics imposes on the Aether's microscopic
structure and shows that they are fully consistent with---and enrich---every
previously derived result.


\subsection{The Quantum Superfluid Requirement}

The spin-statistics connection derived from vortex topology requires three
properties of the medium that go beyond a classical fluid:

\begin{enumerate}[leftmargin=2em]
  \item \textbf{Quantized circulation.}\ The electron spin vortex carries angular
    momentum $\pm\hbar/2$ and exhibits $2\pi$ phase winding.\ This requires the
    circulation to be quantized: the line integral of the velocity field around
    any vortex core takes discrete values $\kappa = n\kappa_0$, where $\kappa_0$
    is the circulation quantum.\ In a superfluid, $\kappa_0 = h/m_A$, where $m_A$
    is the mass of the medium's constituent quantum \cite{donnelly1991}.\ This introduces a
    new fundamental quantity---the Aether quantum mass $m_A$.

  \item \textbf{Spinor order parameter.}\ The existence of two distinct vortex
    orientations (spin up and spin down) that are related by the SU(2) double
    cover of SO(3) requires the condensate to be described by a multi-component
    order parameter $\Psi = (\psi_\uparrow, \psi_\downarrow)$, not a simple
    scalar field.\ This spinor structure is what makes the belt trick
    physical---the two components provide the internal degree of freedom that
    distinguishes the two circulation directions while preserving the topological
    properties needed for the exchange phase $(-1)^{2S}$.

  \item \textbf{Bosonic field operators.}\ The medium itself must obey bosonic
    commutation relations $[\hat{a}, \hat{a}^\dagger] = 1$, permitting unlimited
    occupation of each mode.\ This is consistent with electromagnetic waves
    propagating as bosonic quasiparticles (photons as phonon-like excitations of
    the condensate), with the linear superposition of waves, and with the
    transient binding mechanism in which vacuum fluctuations produce
    particle--antiparticle pairs from the bosonic ground state.
\end{enumerate}

These three requirements are not new postulates---they are logical consequences
of the vortex topology that already appears in the theory.\ The Pauli exclusion
principle cannot be derived from a medium that lacks any one of them.


\subsection{The Equation of State: Why the Aether Is Not a Simple BEC}

The theory's most experimentally constrained property is the invariance of the
bulk modulus $\Ka$.\ The Fresnel drag coefficient confirms $\Ka$ is constant at
density ratios up to $n^2 \approx 6$ (diamond), and the Lorentz factor's
precision to $10^{-8}$ confirms it across all tested velocity regimes.

For a standard weakly interacting Bose--Einstein condensate (BEC), the bulk
modulus scales as $K_{\mathrm{BEC}} \propto n^2$ (where $n$ is the number
density), because the pressure goes as $P = g n^2/2$ with constant interaction
strength $g$.\ This means a simple BEC would have a density-dependent bulk
modulus---contradicting the confirmed invariance of $\Ka$.

The requirement $\Ka = \text{constant}$ constrains the equation of state.\ Since
$\Ka = \rhoa\,(\partial P / \partial \rhoa)$, a constant bulk modulus implies:
%
\begin{equation}
  \frac{\partial P}{\partial \rhoa} = \frac{\Ka}{\rhoa}
  \qquad\Longrightarrow\qquad
  P = \Ka \ln\!\left(\frac{\rhoa}{\rho_0}\right) + P_0
\end{equation}
%
where $\rho_0$ is the reference (ambient far-field) density of the Aether and
$P_0$ is the corresponding ambient pressure.\ Throughout this paper, $\rhoa$
denotes the local Aether density at an arbitrary point, which may differ from
$\rho_0$ due to gravitational compression or cosmological variation.

This \emph{logarithmic equation of state} arises in quantum field theories with
logarithmic nonlinearities \cite{rosen1969,zloshchastiev2011} and has been studied in the context of
superfluids, nuclear matter, and Bose liquids.\ The Aether's equation of state
is therefore not that of a dilute gas of weakly interacting bosons, but of a
denser quantum liquid whose self-interaction potential produces constant
stiffness across all density regimes.

The logarithmic nonlinearity $g(n) = b\,\ln(n/n_0)$ studied in the context
of quantum condensates \cite{rosen1969,zloshchastiev2011} produces a
density-independent speed of sound (from $P = bn$, a linear pressure law).\
This differs from the Aether's logarithmic EOS ($P = \Ka\ln(\rho/\rho_0)$,
constant bulk modulus, density-\emph{dependent} sound speed), though both
are called ``logarithmic.''\ The microscopic interaction that produces
the Aether's EOS can be derived from the Gibbs--Duhem relation
$dP = n\,d\mu$ with $P = \Ka\ln(n/n_0)$, which yields the chemical
potential
%
\begin{equation}\label{eq:gp-nonlinearity}
  \mu(n) = \mu_0 + \Ka\!\left(\frac{1}{n_0} - \frac{1}{n}\right)
\end{equation}
%
where $\mu_0 = \mu(n_0)$ is the equilibrium chemical potential.\ In the
Gross--Pitaevskii equation, $g(n) = \mu(n)$, so the effective nonlinearity
is an \emph{inverse-density} interaction ($-\Ka/n$ plus a constant)---repulsive
at low density (preventing cavitation) and asymptotically constant at high
density.\ This produces exactly $K = \Ka$ (constant) and
$c_s = \sqrt{\Ka/\rho}$ (density-dependent).\
It is distinct from the logarithmic \emph{potential} $b\ln(n)$, which
yields a logarithmic \emph{chemical potential} but a linear pressure.\
The physical origin of the $\Ka/n$ interaction---whether from
effective multi-body scattering, quantum corrections to the condensate
self-energy, or the Aether's internal structure---remains an open question,
but the functional form itself is uniquely determined by the
$\Ka$-invariance constraint.


\subsection{The Healing Length and the UV Cutoff}

In any quantum superfluid, vortex cores have a characteristic size---the
\emph{healing length}---below which the condensate density drops to zero:
%
\begin{equation}
  \xi = \frac{\hbar}{m_A\, c\, \sqrt{2}}
\end{equation}
%
where $c = \sqrt{\Ka/\rhoa}$ is the local wave speed.\ The healing length sets the
minimum scale at which the continuum description of the medium remains valid.

This quantity has a natural identification in the existing theory.\ The Lamb
shift calculation uses the Compton wavelength $\lambdabar_C = \hbar/(m_e c) =
386$~fm as the ultraviolet cutoff for vacuum fluctuations.\ If $m_A = m_e$:
%
\begin{equation}
  \xi = \frac{\lambdabar_C}{\sqrt{2}} \approx 273\;\text{fm}
\end{equation}
%
The Compton wavelength and the healing length are the same order of magnitude.\
What was previously an \emph{ad hoc} cutoff in the QED calculation now has a
physical interpretation: it is the scale at which the internal structure of the
Aether vortex core becomes relevant, and the continuum medium description breaks
down.\ Below $\xi$, the discrete quantum structure of the Aether dominates.

If $m_A \gg m_e$, then $\xi \ll \lambdabar_C$, and the vortex core is much
smaller than the Compton wavelength.\ In this case, the Compton wavelength would
represent a larger-scale feature of the vortex (such as the loop radius $r_{\mathrm{loop}}$)
rather than the core size.


\subsection{Electron Mass from Vortex Energy}

The energy of a quantized vortex in a superfluid is determined entirely by the
medium properties and the vortex geometry.\ For a vortex ring of radius
$r_{\mathrm{loop}}$ with core size $\xi$, the kinetic energy stored in the
circulatory flow is \cite{donnelly1991}:
%
\begin{equation}\label{eq:vortex-energy}
  E_{\mathrm{vortex}} = \frac{1}{2}\,\rhoa\,\kappa_e^2\,r_{\mathrm{loop}}
    \left[\ln\!\left(\frac{8\,r_{\mathrm{loop}}}{\xi}\right) - \frac{1}{2}\right]
\end{equation}
%
where $\kappa_e = h/(2m_A)$ is the circulation quantum for a spin-$\tfrac{1}{2}$
vortex (half the fundamental quantum $\kappa_0 = h/m_A$).

If the electron \emph{is} a vortex ring, its rest energy must equal this kinetic
energy:
%
\begin{equation}\label{eq:mass-from-topology}
  m_e\,c^2 = \frac{1}{2}\,\rhoa\,\kappa_e^2\,r_{\mathrm{loop}}
    \left[\ln\!\left(\frac{8\,r_{\mathrm{loop}}}{\xi}\right) - \frac{1}{2}\right]
\end{equation}
%
This is a remarkable result: \emph{the electron mass is not a free parameter but
is determined by the Aether density, the circulation quantum, and the vortex
geometry}.\ Mass arises from topology---the kinetic energy locked into a
persistent circulatory flow pattern in the medium.

Substituting $\kappa_e = h/(2m_A)$ and $\xi = \hbar/(m_A c \sqrt{2})$:
%
\begin{equation}
  m_e = \frac{\pi^2 \hbar^2\,\rhoa\,r_{\mathrm{loop}}}{2\,m_A^2\,c^2}
    \left[\ln\!\left(\frac{8\,r_{\mathrm{loop}}\,m_A\,c\,\sqrt{2}}{\hbar}\right)
    - \frac{1}{2}\right]
\end{equation}
%
This equation connects three previously independent quantities---$m_e$, $\rhoa$,
and $m_A$---through the vortex geometry.\ Combined with $c^2 = \Ka/\rhoa$, it
provides a new constraint equation that partially breaks the $\rhoa$ degeneracy
identified in the force constant analysis.


\subsection{Vortex Geometry and Length Contraction}

The vortex energy equation (\ref{eq:mass-from-topology}) has an important
consequence for the spatial variation of vortex size.\ The Aether inflow velocity
varies with location (strong near matter, weak in voids), and the two-fluid
superfluid fraction varies with temperature.\ Since $m_e c^2$ must remain
invariant (the electron mass is observed to be the same everywhere), the vortex
geometry must adjust to compensate for the local flow environment.

Rewriting equation~(\ref{eq:mass-from-topology}) using $c^2 = \Ka/\rhoa$:
%
\begin{equation}
  m_e \Ka = \frac{1}{2}\,\rhoa^2\,\kappa_e^2\,r_{\mathrm{loop}}
    \left[\ln\!\left(\frac{8\,r_{\mathrm{loop}}}{\xi}\right) - \frac{1}{2}\right]
\end{equation}
%
Since $\Ka$, $m_e$, and $\kappa_e$ are all invariant, the right-hand side must
be constant:
%
\begin{equation}
  r_{\mathrm{loop}} \left[\ln\!\left(\frac{8\,r_{\mathrm{loop}}}{\xi}\right)
  - \frac{1}{2}\right] \propto \frac{1}{\rhoa^2}
\end{equation}
%
Because $\xi = \hbar/(m_A c\sqrt{2})$ itself depends on density through
$c = \sqrt{\Ka/\rhoa}$, the healing length grows as $\xi \propto \sqrt{\rhoa}$
in denser regions.\ The logarithmic factor $\ln(8r_{\mathrm{loop}}/\xi) - 1/2$
therefore varies weakly (logarithmically) with density, so the dominant
scaling remains $r_{\mathrm{loop}} \propto 1/\rhoa^2$ to leading order.

In regions of effectively higher Aether density, the vortex ring radius shrinks.\
This provides a microscopic \emph{consistency check} on length contraction.\
The macroscopic contraction itself is derived in
Paper~2~\cite{otto-relativity}: the wave equation governing all binding
forces is Lorentz-invariant \cite{lorentz1904}, so any rest-frame
equilibrium maps to a moving equilibrium contracted by $1/\gamma$ in the
direction of motion.

At the vortex level, when an electron moves through the Aether, the dynamic
compression ahead of it (Bernoulli: $\delta\rho/\rhoa \approx v^2/(2c^2)$)
creates a denser rest-frame environment, and the vortex adjusts its geometry to
maintain $m_e c^2 = E_{\mathrm{vortex}}$ in the compressed medium.\ This
dynamic compression is physically distinct from gravitational inflow, where
Bernoulli gives $\Delta\rho = 0$ exactly (free-fall kinetic energy balances the
gravitational potential).\ A moving object does real work on the medium,
creating a genuine density increase.

The vortex scaling $r_{\mathrm{loop}} \propto 1/\rhoa^2$ predicts a contraction
proportional to $1/\gamma^2$ at leading order, while the exact geometric
result is $1/\gamma$.\ The factor-of-two discrepancy in the exponent arises
because the Donnelly vortex energy formula is non-relativistic---it assumes a
vortex at rest in the fluid.\ Additionally, the Bernoulli stagnation pressure
$v^2/(2c^2)$ is the \emph{maximum} compression (directly ahead of the object);
the actual density field is angle-dependent, and the effective compression
experienced by a point vortex in a flowing medium differs from the uniform
stagnation value.\ A relativistic treatment of vortex dynamics in a flowing
superfluid would be needed to close this gap (see Open Questions).


\subsection{The Aether Quantum Mass}

The circulation quantum $\kappa_0 = h/m_A$ introduces a new fundamental
quantity: the mass $m_A$ of the Aether's constituent quantum.\ This is the mass
scale that determines the vortex core size, the circulation quantum, and---through
the vortex energy equation---the relationship between particle mass and Aether
density.

The linearity constraint provides an important bound on $m_A$.\ The vacuum
fluctuation amplitude satisfies:
%
\begin{equation}
  \left\langle\left(\frac{\delta\rho}{\rhoa}\right)^{\!2}\right\rangle
  = \frac{2\,u_{\mathrm{vac}}}{\rhoa\,c^2}
  \approx \frac{2 \times 10^5}{\rhoa}\;\text{kg/m}^3
\end{equation}
%
For the linear approximation to hold ($\delta\rho/\rhoa \ll 1$), we need
$\rhoa \gg 2 \times 10^5$~kg/m$^3$.\ The vortex energy equation then constrains
$m_A$:

\begin{itemize}[leftmargin=2em]
  \item If $m_A = m_e$: the vortex energy equation gives $\rhoa \sim 10^6$~kg/m$^3$,
    which is marginally linear ($\delta\rho/\rhoa \sim 0.45$).\ The confirmed
    precision of the Lorentz factor to $10^{-8}$ makes this uncomfortable.

  \item If $m_A \sim 100\,m_e$: the equation gives $\rhoa \sim 10^{10}$~kg/m$^3$,
    with $\delta\rho/\rhoa \sim 0.005$---comfortably within the linear regime.

  \item If $m_A \sim m_{\mathrm{Planck}}$: $\rhoa$ becomes extremely large and the
    medium is extraordinarily stiff ($\Ka \sim 10^{46}$~J/m$^3$), with vacuum
    fluctuations completely negligible relative to the background density.
\end{itemize}

The experimental data thus \emph{favors} $m_A \gg m_e$: the heavier the Aether
quantum, the higher $\rhoa$, and the more comfortably the linearity constraint
is satisfied.\ This is consistent with the gamma-ray timing constraints from the
Fermi Large Area Telescope \cite{vasileiou2013}, which constrain Lorentz-violating
photon dispersion to energy scales above $7.6\,E_{\mathrm{Planck}}$.\ Since
photons propagate as spin-wave excitations (microscopically distinct from
density waves, though macroscopically equivalent through spin-orbit
coupling), this constrains the spin healing length $\xi_s$, not the
density healing length $\xi$---consistent
with the spin channel being far stiffer than the density channel.



\subsection{Topological Stability and Derrick's Theorem}
\label{sec:derrick}

The quantum superfluid framework raises an essential question: could the electron be
a \emph{hedgehog skyrmion}---a spherically symmetric texture in the spinor order
parameter---rather than a quantized vortex ring?\ Skyrmions are topologically
nontrivial field configurations classified by the homotopy group $\pi_3(S^3) = \mathbb{Z}$,
and they serve as particle-like solitons in nuclear physics (the Skyrme model
\cite{skyrme1962}).\ In three spatial dimensions, the hedgehog ansatz takes the form
$\hat{n}(\mathbf{r}) = \hat{r}$ with a radial profile $F(r)$ satisfying $F(0) = \pi$
and $F(\infty) = 0$, wrapping the order parameter sphere exactly once.

However, Derrick's theorem \cite{derrick1964} places a fatal constraint on such
configurations in the logarithmic condensate.\ Under a uniform spatial rescaling
$\psi(\mathbf{r}) \to \psi(\mathbf{r}/L)$, the energy functional separates into
terms with definite scaling behavior:
%
\begin{equation}
  E_{\mathrm{gradient}} \propto L, \qquad
  E_{\mathrm{potential}} \propto L^3
\end{equation}
%
The gradient energy (both density and spin contributions) scales as $L$ because
$|\nabla\psi|^2$ contributes a factor $1/L^2$ while the volume element contributes
$L^3$.\ The logarithmic potential energy $V(\psi) = \Ka[\,|\psi|^2(2\ln|\psi| - 1) + 1\,]$
scales as $L^3$ from the volume element alone.\ The total energy is therefore:
%
\begin{equation}
  E(L) = A\,L + B\,L^3 \qquad (A, B > 0)
\end{equation}
%
Since both coefficients are strictly positive, $dE/dL = A + 3BL^2 > 0$ for all $L > 0$.
The energy is \emph{monotonically increasing}---the configuration always lowers its
energy by shrinking ($L \to 0$).\ There is no stable skyrmion.

Numerical verification confirms this conclusion.\ Using the coupled Euler--Lagrange
equations for the hedgehog energy functional
%
\begin{equation}
  E = C \cdot \Ka \cdot \xi^3 = 4\pi\Ka\xi^3 \int_0^\infty s^2\,ds\;
    \bigl[\,g'^2 + g^2(F'^2 + \sin^2\!F/s^2) + g^2(2\ln g - 1) + 1\,\bigr]
\end{equation}
%
where $s = r/\xi$ and $g = |\psi|/\psi_\infty$, iterative relaxation from physically
reasonable initial profiles drives the spin texture $F(s)$ toward zero everywhere
and the dimensionless energy coefficient $C$ toward zero---the skyrmion collapses
to a point.

This stands in sharp contrast to the Skyrme model, where a quartic derivative
term $\propto |\nabla\psi|^4$ contributes an energy that scales as $1/L$:
%
\begin{equation}
  E_{\mathrm{Skyrme}}(L) = A\,L + \frac{D}{L} \qquad \Rightarrow \qquad
  L_{\mathrm{min}} = \sqrt{D/A}
\end{equation}
%
The competition between the $L$ and $1/L$ terms creates a stable minimum at finite
size.\ The logarithmic condensate has no such quartic term---its equation of state
is fundamentally softer.

This result has a profound physical consequence: \textbf{the electron cannot be a
hedgehog skyrmion in the logarithmic condensate.}\ It must instead be a quantized
vortex ring, whose stability rests on an entirely different mechanism---circulation
quantization:
%
\begin{equation}
  \kappa = n\,\frac{h}{m_A} \qquad (n \in \mathbb{Z})
\end{equation}
%
The quantized circulation $\kappa$ is a \emph{discrete} topological invariant.\ Unlike
the continuous hedgehog texture, which can be smoothly deformed to zero (Derrick
collapse), a vortex ring's winding number is exactly integer and cannot change
continuously.\ The vortex core size is fixed at the healing length
$\xi = \hbar/(m_A c\sqrt{2})$, and the circulation quantum $\kappa_0 = h/m_A$ is
determined by the microscopic structure of the medium.\ Neither can be reduced
by any smooth deformation of the field.

The distinction is between two types of topological protection:
%
\begin{itemize}[leftmargin=2em]
  \item \textbf{Homotopy (skyrmion):} $\pi_3(S^3) = \mathbb{Z}$.\ The winding number
    is a topological invariant under continuous deformations that preserve the boundary
    conditions.\ However, Derrick's theorem shows that the energetically favored
    deformation is $L \to 0$---the configuration \emph{collapses} rather than unwinding.
    The topological charge is conserved but concentrated into a singular point.

  \item \textbf{Circulation quantization (vortex):} $\kappa = n\,h/m_A$.\ The
    quantization is enforced by the single-valuedness of the condensate wave function
    around any closed loop encircling the vortex core.\ This is a \emph{stronger}
    constraint: the core cannot collapse because $\xi$ is fixed by the healing length,
    and the circulation cannot decrease because $\kappa_0$ is fixed by $h/m_A$.
\end{itemize}
%
The Derrick instability of the hedgehog skyrmion therefore provides independent
theoretical support for identifying the electron as a quantized vortex ring rather
than a field-theoretic soliton.\ The electron's stability is not an energetic
balance between competing gradient terms (as in the Skyrme model) but a consequence
of discrete topological quantization inherent to the superfluid medium.



\subsection{Constraints from the Strong Force}
\label{sec:strong-constraints}

The identification of the strong force as a compressed-Aether bridge
(Section~\ref{sec:strong-force}) provides an independent set of constraints on
the Aether's microscopic properties.\ The QCD string tension, the deconfinement
temperature, and the flux tube geometry are all precisely measured, and each
imposes a quantitative requirement on the same parameters---$\rhoa$, $\Ka$,
$m_A$, and $\xi$---that appear in the vortex energy equation and the linearity
constraint.

\subsubsection*{Bridge diameter and the healing length}

The deconfinement temperature $T_{\mathrm{deconf}} \approx 170$~MeV determines
the bridge diameter through the thermal energy balance: the bridge dissolves when
the thermal energy per unit cross-section equals the string tension:
%
\begin{equation}
  d_{\mathrm{bridge}} \approx \frac{k_B T_{\mathrm{deconf}}}{\sigma}
    \approx \frac{170\;\mathrm{MeV}}{1\;\mathrm{GeV/fm}} \approx 0.17\;\mathrm{fm}
  \label{eq:bridge-diameter}
\end{equation}
%
Lattice QCD simulations measure the flux tube width (FWHM of the chromoelectric
field distribution) at approximately 0.5~fm \cite{cardoso2013,cea2024}, somewhat
larger than this core estimate because the field profile has tails beyond the core.

The bridge is a condensate structure, and the minimum scale of any structure in
a quantum superfluid is the healing length $\xi = \hbar/(m_A\,c_{\mathrm{local}}\sqrt{2})$,
which decreases inside the compressed bridge where $c_{\mathrm{local}} < c$,
making the bridge core even more compact.\ The bridge core diameter must
therefore satisfy $d_{\mathrm{bridge}} \gtrsim \xi$, giving:
%
\begin{equation}
  \xi \lesssim 0.17\;\mathrm{fm}
  \label{eq:xi-upper-bound}
\end{equation}
%
Combined with the definition $\xi = \hbar / (m_A c \sqrt{2})$, this constrains
the Aether quantum mass:
%
\begin{equation}
  m_A \gtrsim \frac{\hbar}{\xi_{\max}\, c\, \sqrt{2}}
    \approx \frac{1.055 \times 10^{-34}}{1.7 \times 10^{-16} \times 3 \times 10^8
    \times 1.414} \approx 1.5 \times 10^{-27}\;\mathrm{kg}
    \approx 820\;\mathrm{MeV}/c^2
  \label{eq:mA-lower-bound}
\end{equation}
%
If the bridge diameter is several times $\xi$ (as expected, since the
condensation process typically extends over a few healing lengths), the
constraint relaxes.\ However, three independent lines of evidence converge on
a substantially narrower range:
%
\begin{enumerate}[leftmargin=2em]
  \item \textbf{Bridge diameter from lattice QCD.}\ High-precision lattice
    simulations measure the chromoelectric flux tube core width at
    $d_{\mathrm{core}} \approx 0.30$--$0.35$~fm \cite{cardoso2013}.\ If the
    bridge diameter equals $d_{\mathrm{core}}$ and
    $d_{\mathrm{bridge}} \sim 1$--$2\,\xi$, then
    $\xi \approx 0.15$--$0.35$~fm, giving
    $m_A \approx 470$--$930$~MeV/$c^2$.

  \item \textbf{Self-consistent vortex--bridge equations.}\ Simultaneously
    solving the electron vortex energy equation~\eqref{eq:mass-from-topology}
    and the bridge compression ratio with the lattice-measured string tension
    $\sigma \approx 0.18$~GeV$^2$ yields a self-consistent solution in the
    range $m_A \approx 600$--$1000$~MeV/$c^2$.

  \item \textbf{Glueball mass spectrum.}\ If Aether quanta are the
    fundamental degrees of freedom, the lightest glueball (a bound state of
    pure medium excitations) provides an upper bound on $m_A$.\ Lattice QCD
    finds $m_{0^{++}} \approx 1710$~MeV/$c^2$ for the lightest scalar
    glueball \cite{morningstar1999}, but the lowest-mass glueball candidate
    consistent with experimental data sits near $\sim$1500--1700~MeV.\ Since
    a glueball must contain at least two quanta, this suggests
    $m_A \lesssim 850$~MeV/$c^2$.
\end{enumerate}
%
The overlap of these three constraints yields a \emph{convergence zone}:
%
\begin{equation}
  500\;\mathrm{MeV}/c^2 \;\lesssim\; m_A \;\lesssim\; 850\;\mathrm{MeV}/c^2
  \label{eq:mA-range}
\end{equation}
%
This range corresponds to $\xi \approx 0.17$--$0.28$~fm,
$\rhoa \approx 4 \times 10^{11}$--$10^{12}$~kg/m$^3$, and
$\Ka \approx 4 \times 10^{28}$--$10^{29}$~J/m$^3$---all consistent with the
values derived independently from the electron vortex equation and the bridge
compression ratio (Table above).\ The convergence of $m_A$ to the hadronic
mass scale is physically natural: the Aether quantum is the fundamental unit
of the medium, and hadrons are its primary bound structures.


\subsubsection*{Two density scales: ambient and bound}

The string tension $\sigma \approx 1$~GeV/fm imposes a constraint on the energy
density inside the bridge.\ If the bridge has cross-sectional area
$A \approx \pi\xi^2/4$ (taking the bridge diameter as approximately $\xi$ in the minimum-size scenario), the internal energy density is:
%
\begin{equation}
  \varepsilon_{\mathrm{bridge}} = \frac{\sigma}{A}
    \approx \frac{1.6 \times 10^5\;\mathrm{J/m}}{\pi\,\xi^2/4}
    \sim 10^{36}\text{--}10^{37}\;\mathrm{J/m^3}
  \label{eq:bridge-energy-density}
\end{equation}
%
If this energy density equaled $\Ka$, it would imply
$\rhoa = \Ka/c^2 \sim 10^{19}\text{--}10^{20}$~kg/m$^3$.\ However, the
vortex energy equation for the electron
%
\begin{equation}
  m_e c^2 = \tfrac{1}{2}\,\rhoa\,\kappa_e^2\,r_{\mathrm{loop}}\,
    [\ln(8r_{\mathrm{loop}}/\xi) - \tfrac{1}{2}]
\end{equation}
%
with $r_{\mathrm{loop}} \approx \lambdabar_C \approx 386$~fm requires a much
lower ambient density:
%
\begin{equation}
  \rhoa \sim 10^{11}\text{--}10^{12}\;\mathrm{kg/m^3}, \qquad
  \Ka = \rhoa\,c^2 \sim 10^{28}\text{--}10^{29}\;\mathrm{J/m^3}
  \label{eq:ambient-density}
\end{equation}
%
The two constraints are not contradictory---they reveal that \textbf{the binding
process compresses Aether by a factor of approximately $7 \times 10^7$}.\ The ambient
superfluid has density $\rhoa \sim 10^{11}$~kg/m$^3$ and bulk modulus
$\Ka \sim 10^{28}$~J/m$^3$.\ The compressed material inside a bridge has an effective
energy density $\sim 10^{36}$~J/m$^3$---approximately $7 \times 10^7$ times the ambient
$\Ka$.

This compression ratio is not a free parameter---it is algebraically determined
by three measured or constrained quantities.\ The bridge energy density is
$\varepsilon_{\mathrm{bridge}} = \sigma / A = 4\sigma / (\pi\xi^2)$, and the
compression ratio is:
%
\begin{equation}
  X \equiv \frac{\rho_{\mathrm{bridge}}}{\rhoa}
    = \frac{\varepsilon_{\mathrm{bridge}}}{\Ka}
    = \frac{4\sigma}{\pi\,\xi^2\,\Ka}
  \label{eq:compression-ratio}
\end{equation}
%
The logarithmic equation of state $P = \Ka \ln(\rho/\rho_0) + P_0$---here
extrapolated well beyond the experimentally tested regime, assuming $\Ka$ remains
constant at extreme compressions (see Open Question~10)---governs
how much work is required to achieve this compression.\ For a parcel of Aether
compressed from $\rho_0$ to $\rho_f = X \rho_0$, the energy per unit volume of
the compressed material is:
%
\begin{equation}
  \varepsilon_{\mathrm{compress}} = \Ka\bigl[X - \ln X - 1\bigr]
  \label{eq:compression-energy}
\end{equation}
%
For large $X$, $\ln X \ll X$, so $\varepsilon_{\mathrm{compress}} \approx
\Ka\, X = \varepsilon_{\mathrm{bridge}}$: the compression work accounts for
essentially all of the bridge energy density.\ The bridge is a compressed
strand of Aether whose energy is almost entirely stored as work against the
medium's bulk modulus.\ The logarithmic
term $\ln X \approx 18$ (for $X \approx 7 \times 10^7$) provides a small correction
of order $10^{-7}$ relative to $X$.

This clarifies the relationship between the ``bound'' and ``compressed'' descriptions
of bridge material used in earlier sections.\ The bridge is not a phase transition
in the thermodynamic sense (unlike the electron vortex, whose stability rests on
topological circulation quantization).\ Rather, it is a region of extreme elastic
compression \emph{stabilized by the quark vortex endpoints that anchor it}.\ The
quark vortices at each end provide the topological boundary conditions that maintain
the compressed strand; without them, the compression would immediately relax back
to ambient density.\ The deconfinement ``melting'' described in
Section~\ref{sec:strong-force} is the thermal disruption of these anchoring vortex
endpoints, which releases the elastic energy and allows the compressed Aether to
expand back to its ambient state.\ The bridge material thus remains in the same
thermodynamic phase as the surrounding superfluid---it is simply compressed to
$\sim\!7 \times 10^7$ times ambient density by the topological constraints at its
endpoints.

\paragraph{Cross-check from the electron vortex equation.}
Using the electron-derived $\Ka$ values from
equation~\eqref{eq:ambient-density} with a range of healing-length assumptions
(including values outside the narrowed range $\xi = 0.17$--$0.28$~fm to test
robustness) yields a remarkably consistent compression ratio:
%
\begin{center}
\begin{tabular}{@{}lccc@{}}
\hline
\textbf{Scenario} & $\Ka$ (J/m$^3$) & $X = 4\sigma/(\pi\xi^2\Ka)$ & $\log_{10} X$ \\
\hline
$\xi = 0.17$~fm  & $9.6 \times 10^{28}$ & $7.3 \times 10^{7}$ & 7.87 \\
$\xi = 0.10$~fm  & $2.6 \times 10^{29}$ & $7.8 \times 10^{7}$ & 7.89 \\
$\xi = 0.06$~fm  & $6.8 \times 10^{29}$ & $8.3 \times 10^{7}$ & 7.92 \\
$\xi = 0.40$~fm  & $2.0 \times 10^{28}$ & $6.5 \times 10^{7}$ & 7.81 \\
\hline
\end{tabular}
\end{center}
%
Across all scenarios, $\log_{10} X = 7.85 \pm 0.05$---the compression ratio is
pinned at $X \approx 7 \times 10^7$, independent of which $\xi$ is assumed.\
This convergence occurs because both $\Ka$ and $\xi$ are determined by the same
underlying parameters ($\rhoa$, $m_A$, $c$), so their $\xi$-dependence cancels
in the ratio.

\paragraph{Physical consequences of the compression.}
The compressed bridge interior has qualitatively different properties from the
ambient superfluid.\ The speed of sound in the logarithmic condensate is
$c_s = \sqrt{\Ka/\rho}$; in the bridge, $c_{s,\mathrm{bridge}} = c / \sqrt{X}
\approx 10^{-4}\,c \approx 3 \times 10^4$~m/s.\ Excitations inside the bridge
propagate $10{,}000$ times slower than in the ambient medium---the bridge
interior is a slow, dense conduit, qualitatively consistent with the behavior
of gluons versus photons (which propagate
freely through the ambient superfluid).
The slow propagation speed explains the mass scale and self-interaction
properties of gluons; confinement itself arises from the bridge-snapping
mechanism described in Section~\ref{sec:strong-force}.

\paragraph{Bridge profile and minimum width.}
The bridge diameter $d_{\mathrm{bridge}} \approx 0.30$--$0.35$~fm is only
$1$--$2$ times the healing length $\xi \approx 0.17$--$0.28$~fm.\ A structure
this compact has no extended bulk interior: the entire cross-section lies within
the $\sim\!\xi$-wide transition region where the condensate density rises from
ambient to its peak and returns.\ The compression ratio $X \approx 7 \times 10^7$
therefore represents the \emph{peak} central density, not a uniform interior
value; the volume-averaged compression is lower.

This compactness has two important consequences.\ First, the bridge sits at the
\emph{minimum possible width}: the condensate cannot sustain a compressed
structure narrower than $\sim\!\xi$, so the bridge is as thin as the medium
permits.\ This is energetically optimal---for a fixed string tension $\sigma$,
minimizing the cross-sectional area minimizes the total compressed volume and
hence the total energy per unit length.\ Second, because the bridge profile is
entirely governed by the healing-length recovery dynamics, its radial density
distribution is predicted to follow the Gross--Pitaevskii condensate profile
for a cylindrical perturbation---a Gaussian-like ``Gausson'' characteristic of
logarithmic nonlinearities.\ Lattice QCD simulations measure the transverse
chromoelectric field profile with a width (FWHM) of $w \approx 0.5$~fm at
separations $d \gtrsim 0.7$~fm \cite{cardoso2013,cea2024}, consistent with
the predicted $\sim\!2$--$3\,\xi \approx 0.4$--$0.6$~fm.\ The measured
profile is well fit by both Gaussian and $\operatorname{sech}^2$ functional
forms \cite{cea2024}; current lattice precision cannot distinguish between
these and the GP Gausson, which differ primarily in their tail behaviour.\
Earlier dual-superconductor fits extract a penetration depth
$\lambda \approx 0.22$--$0.24$~fm \cite{cea2012}, overlapping with
$\xi \approx 0.17$--$0.28$~fm.\ The width agreement is non-trivial: the
bridge model predicts it from $\xi$ alone, with no adjustable parameter.\
A future precision measurement of the tail profile---Gaussian versus
$\operatorname{sech}^2$ versus Bessel---would provide a sharper test.

The bridge draws its compressed material from the surrounding volume.\ A
minimum bridge segment (carrying the pair-creation energy $2m_e c^2$) has
length $L_{\min} = 2m_e c^2/\sigma \approx 10^{-3}$~fm and draws Aether from a
sphere of radius $\sim$16~fm, roughly the size of a large nucleus.\ For
quark--antiquark pair creation (hadronization), the bridge must store
enough energy to create a new quark--antiquark pair including
their kinetic energy within the resulting mesons ($\sim$1~GeV total),
requiring a length $L_{\mathrm{snap}} \approx 1.2$~fm---matching the
string-breaking distance measured in lattice QCD simulations
\cite{bali2005,bulava2019}.


\subsubsection*{Refined parameter estimates}

Combining the strong force constraints with the vortex energy equation and the
speed of light yields a consistent set of Aether properties:
%
\begin{center}
\begin{tabular}{@{}lll@{}}
\hline
\textbf{Property} & \textbf{Value} & \textbf{Determined by} \\
\hline
Healing length $\xi$       & $0.17$--$0.28$~fm     & Bridge diameter / lattice QCD \\
Aether quantum mass $m_A$  & $500$--$850$~MeV/$c^2$  & Flux tube width, glueball mass \\
Ambient density $\rhoa$    & $\sim 10^{11}$--$10^{12}$~kg/m$^3$ & Electron vortex energy \\
Bulk modulus $\Ka$          & $\sim 10^{28}$--$10^{29}$~J/m$^3$  & $\Ka = \rhoa\,c^2$ \\
Compression ratio          & $\sim 7 \times 10^{7}$         & $\sigma / (\Ka \cdot \pi\xi^2/4)$ \\
Bridge density              & $\sim 10^{19}$--$10^{20}$~kg/m$^3$ & String tension \\
\hline
\end{tabular}
\end{center}
%
Several cross-checks confirm internal consistency:
%
\begin{itemize}[leftmargin=2em]
  \item \textbf{Linearity.}\ With $\rhoa \sim 10^{11}$~kg/m$^3$, the density
    perturbation from an electron at its Compton wavelength is
    $\delta\rho/\rhoa = 2Gm_e/(\lambdabar_C c^2) \approx 3.5 \times 10^{-45}$,
    confirming the Lorentz factor to extraordinary precision.

  \item \textbf{Proton mass.}\ Three bridge arms of length $\sim$0.3~fm at
    string tension $\sigma \approx 1$~GeV/fm give
    $E_{\mathrm{bridge}} \approx 3 \times 0.3 \times 1 = 900$~MeV, compared
    to the measured proton mass of 938~MeV---the remaining $\sim$38~MeV is
    accounted for by the quark endpoint masses ($\sim$9~MeV) and the junction
    binding energy.

  \item \textbf{Fermi LAT and the spin-wave cutoff.}\ The density healing
    length $\xi \approx 0.17$--$0.28$~fm corresponds to a UV cutoff energy
    $E_{\mathrm{UV}} = \hbar c / \xi \approx 700$--$1200$~MeV.\ If photons
    were density-wave (Bogoliubov) excitations, their dispersion would
    deviate from $\omega = ck$ at this scale---inconsistent with the Fermi
    LAT quadratic constraint ($E_{\mathrm{QG},2} \gtrsim 10^{11}$~GeV).\
    However, the theory identifies photons as \emph{spin-wave} excitations
    of the spinor condensate, which have a separate spin healing length
    $\xi_s = \hbar / \sqrt{2\,m_A\,q}$, where $q$ is the spin-dependent
    coupling strength.\ The Fermi LAT constraint requires
    $\xi_s \ll \xi$, i.e., $q \gg \tilde{g}\,n_0$.\ This is physically
    consistent: the electromagnetic force (spin-wave channel) is
    $\sim\!10^{36}$ times stronger than gravity (density channel), so
    the spin coupling should vastly exceed the density coupling.\ The
    Fermi LAT thus provides an independent confirmation that the photon
    channel is the \emph{stiff} degree of freedom of the condensate,
    while gravity operates through the much softer density channel.

  \item \textbf{UV cutoff at the hadronic scale.}\ The energy scale
    $E_{\mathrm{UV}} = \hbar c / \xi \approx 700$--$1200$~MeV sits precisely
    at the QCD confinement scale, where perturbative QFT gives way to
    non-perturbative hadron physics.\ This is not a coincidence: both the
    theory's continuum breakdown and the onset of confinement reflect the
    same underlying transition---the scale at which individual Aether quanta
    become resolvable.

  \item \textbf{Bridge compactness.}\ With $d_{\mathrm{bridge}} \approx
    0.30$--$0.35$~fm and $\xi \approx 0.17$--$0.28$~fm, the ratio
    $d_{\mathrm{bridge}}/\xi \approx 1.1$--$2.1$: bridges are only
    1--2 healing lengths across.\ Unlike most superfluid structures, which
    extend over many healing lengths, the bridge is essentially
    ``all core''---compressed to the absolute minimum scale the condensate
    permits.\ This maximal compactness is consistent with the bridge
    carrying the maximum possible energy density for a given cross-section.

  \item \textbf{Vortex ring geometry.}\ The narrowed $\xi$ pins the
    electron's aspect ratio to $r_{\mathrm{loop}}/\xi \approx
    1{,}400$--$2{,}300$, giving $\ln(8r_{\mathrm{loop}}/\xi) - 1/2 \approx
    8.8$--$9.3$ ($\pm 3\%$ from the mean).\ The electron is a macroscopic
    vortex ring relative to the medium's fundamental scale---a large, stable
    structure far above the core size, consistent with its exceptional
    stability.\ The near-constancy of the logarithmic factor means the
    electron mass equation determines $\rhoa$ to better than one significant
    figure once $m_A$ is known.

  \item \textbf{Proton geometry.}\ The proton charge radius
    ($r_p \approx 0.84$~fm) spans only $3$--$5\,\xi$, and each Y-junction
    arm ($\sim$0.3~fm) is barely $1$--$1.8\,\xi$.\ The proton's internal
    structure operates at the absolute minimum scale of the condensate, with
    no extended ``bulk bridge'' region---consistent with the difficulty of
    perturbative QCD at low momentum transfer.

  \item \textbf{String-breaking distance.}\ The hadronization length
    $L_{\mathrm{snap}} \approx 1.2$~fm corresponds to $4$--$7\,\xi$:
    a bridge must stretch several healing lengths beyond equilibrium
    before pair creation becomes energetically favorable.\ Lattice QCD
    simulations measure the string-breaking distance at
    $r_{\mathrm{break}} \approx 1.2$~fm \cite{bali2005,bulava2019},
    in direct agreement with this prediction.
\end{itemize}




\subsection{Emergence of the Schr\"{o}dinger Equation}

The theory's quantum superfluid structure contains all the ingredients needed to derive
the Schr\"{o}dinger equation for quasiparticle excitations.\ This derivation is not a
new postulate---it follows from the condensate dynamics already established.

\subsubsection*{The Madelung representation}

The Gross--Pitaevskii equation for the Aether order parameter
$\psi = \sqrt{\rho/m_A}\,e^{i\phi}$ decomposes via the Madelung
transformation \cite{madelung1927} into continuity and Euler equations,
with a quantum potential
$Q = -(\hbar^2/2m_A)\,\nabla^2\!\sqrt{\rho}/\sqrt{\rho}$ encoding all
interference and diffraction effects.\ The inverse transformation
recovers the nonlinear Schr\"{o}dinger equation (Gross--Pitaevskii), which
for small perturbations around a uniform background linearizes to the
ordinary Schr\"{o}dinger equation for quasiparticle excitations.

\subsubsection*{The Wallstrom objection and its resolution}

Wallstrom \cite{wallstrom1994} raised a well-known objection to hydrodynamic derivations
of quantum mechanics: the Madelung equations are necessary but not sufficient, because
they do not enforce single-valuedness of $\psi$.\ A classical fluid can have arbitrary
circulation, but quantum mechanics requires the phase $\phi$ to change by $2\pi n$ around
any closed loop, giving quantized circulation $\oint\mathbf{v}\cdot d\mathbf{l} = n\,h/m_A$.

In the Aether framework, this objection is resolved automatically.\ The medium is a
quantum superfluid whose order parameter $\psi$ is single-valued by construction---this
is the same property that produces quantized vortices (Section~\ref{sec:qss}).\ The
circulation quantization $\kappa = n\,h/m_A$ is not an additional assumption but is
already built into the theory's fundamental structure.\ The Wallstrom gap---which is
fatal for attempts to derive quantum mechanics from classical fluid mechanics---does not
apply to a genuine quantum superfluid.

\subsubsection*{Bogoliubov linearization}

An independent route---spectral analysis rather than hydrodynamic
transformation---comes from linearizing the Gross--Pitaevskii equation
around the condensate ground state, yielding the Bogoliubov--de~Gennes
equations \cite{bogoliubov1947} for quasiparticle amplitudes.\ The
excitation spectrum transitions from phonon-like ($\omega \propto k$) at
long wavelengths to free-particle-like ($\omega \propto k^2$) at short
wavelengths, with crossover at the healing length $\xi$.\ Both routes
yield the same result.

This provides a concrete physical picture: \textbf{the Schr\"{o}dinger equation describes
the dynamics of quasiparticle excitations above the condensate ground state}.\ Particles
are vortex excitations of the Aether, and their quantum-mechanical behavior---interference,
tunneling, discrete energy levels---is the behavior of quantized excitations in a superfluid,
governed by the linearized condensate dynamics.


% ============================================================

% ============================================================
\section{Summary of Theoretical Properties}
% ============================================================

\subsection{Properties of Light (Paper 7 contribution)}

\begin{itemize}[leftmargin=2em]
  \item Photons are spin-wave excitations of the spinor Aether condensate---transverse
    oscillations of the condensate's internal spin degree of freedom, microscopically
    distinct from longitudinal density compressions.\ At macroscopic wavelengths,
    spin-orbit coupling produces effective mass currents, making photons
    indistinguishable from first-sound perturbations with propagation speed~$c_1$
    (Paper~6~\cite{otto-cosmology}).\ Derrick's theorem forbids stable localized
    particles in the logarithmic condensate, so photons propagate as extended waves
    in free space.\
    Detection occurs when the wave enters the strong-inflow environment near matter,
    where DCE dynamics facilitate condensation into a localized quantum $E = h\nu$.\
    Polarization states correspond to the two spin-wave modes of the spinor order
    parameter.
  \item Through spin-orbit coupling \cite{volovik2003}, these spin-wave excitations
    produce effective mass currents that manifest as the flow-vector (electric) and
    flow-angle (magnetic) fields described in Paper~4~\cite{otto-em}.
\end{itemize}


\subsection{Properties of Atomic Structure}

\begin{enumerate}[leftmargin=2em]
  \item Atoms function as resonant cavities in the Aether, with the nuclear electric field
    (radial Aether inflow) creating a flow gradient that confines electron waves through
    acoustic refraction.
  \item Energy levels are discrete because only resonant modes---those satisfying angular
    phase matching and radial confinement---persist in the cavity.\ Non-resonant modes
    destructively interfere and radiate away.
  \item The quantization condition ($n\lambda = 2\pi r$) is the requirement for
    constructive self-interference of a wave propagating around a closed path in the
    medium.
  \item The ground state energy $E_1 = -\frac{1}{2}\alpha^2 m_e c^2 = -13.6$~eV is
    determined by three quantities with established Aether interpretations: $c$ (wave
    speed), $\alpha$ (electron flow strength), and $m_e c^2$ (binding rest energy).
  \item The fine structure constant $\alpha \approx 1/137$ is the ratio of the electron's
    orbital velocity to the Aether wave speed---an atomic Mach number.
  \item The Bohr radius $a_0 = 0.529$~\AA{} is the characteristic size of the smallest
    resonant mode in the nuclear Aether cavity.
  \item Fine structure arises from the interaction of the electron's spin vortex with
    its orbital motion through the Aether (spin-orbit coupling), producing corrections
    of order $\alpha^2$.
  \item Hyperfine structure arises from the coupling of two nested spin vortices
    (electron and proton), producing the 21~cm hydrogen line.
  \item The Zeeman and Stark effects result from external Aether flows (magnetic
    circulation, electric inflow) distorting the symmetry of the resonant cavity.
  \item The quantum defect in multi-electron atoms measures how deeply the electron's
    resonant mode penetrates the inner flow structure of the Aether cavity.
  \item Energy levels scale as $Z^2$ for hydrogen-like ions, reflecting the stronger
    Aether inflow from higher nuclear charge.
  \item Exotic systems (positronium, muonic hydrogen) follow the same cavity model with
    modified reduced masses, providing independent experimental confirmation.
\end{enumerate}



\subsection{Properties of Fermion Statistics}

\begin{enumerate}[leftmargin=2em]
  \item The Pauli exclusion principle arises from the topological instability of
    multiply quantized vortices in the Aether: two same-circulation vortices
    superpose to create a doubly wound configuration that is energetically unstable
    ($s^2$ vs.\ $s$ energy scaling) and topologically inconsistent.
  \item Each resonant cavity mode admits at most two electron vortices, which must
    have opposite circulation (opposite spin), producing the fundamental
    two-electron-per-orbital rule.
  \item The antisymmetric wavefunction under particle exchange follows from the
    SU(2) topology of spin-$\tfrac{1}{2}$ vortices: exchange is equivalent to a
    $2\pi$ relative rotation, which produces a sign change for half-integer
    winding objects.
  \item Atomic shell capacities ($2n^2$) follow from the resonant cavity mode
    count $\sum_{l=0}^{n-1} 2(2l+1)$ with two opposite-circulation vortices per mode.
  \item Hund's first rule (maximize spin) reflects the minimization of
    vortex--vortex interaction energy: same-circulation vortices in different modes
    have less flow-field overlap than opposite-circulation vortices in the same mode.
  \item Electron degeneracy pressure arises because topological vortex exclusion
    forces electrons into higher-energy cavity modes, creating a pressure that
    resists compression and scales as $n_e^{5/3}$.
  \item The Chandrasekhar limit ($1.44\,M_\odot$) is the point where vortex flow
    speeds in the highest occupied modes approach $c$, the wave speed of the medium.
  \item Integer-spin particles (bosons) have $2\pi$ phase winding per rotation and
    acquire no sign change under exchange, permitting unlimited state occupation---consistent
    with the observed behavior of photons and helium-4 atoms.
  \item Covalent chemical bonds require opposite-circulation vortex pairs in shared
    cavity modes; same-circulation pairs are forbidden by topological exclusion.
  \item The dependence of exchange statistics on spatial dimensionality (fermion/boson
    in 3D, anyons in 2D) confirms the topological origin of the exclusion principle.
\end{enumerate}



\subsection{Properties of Quantum Superfluid Structure}

\begin{enumerate}[leftmargin=2em]
  \item The Aether is a quantum superfluid with quantized circulation
    $\kappa = n \kappa_0$, where $\kappa_0 = h/m_A$ and $m_A$ is the mass of
    the Aether's constituent quantum.
  \item The condensate is described by a spinor (two-component) order parameter
    $\Psi = (\psi_\uparrow, \psi_\downarrow)$, encoding the two possible
    vortex circulations (spin up and spin down).
  \item The Aether field operators obey bosonic commutation relations, consistent
    with linear wave superposition, photon propagation, and vacuum pair
    production.
  \item The equation of state is logarithmic: $P = \Ka \ln(\rhoa/\rho_0) + P_0$,
    the unique form that produces a density-independent bulk modulus $\Ka$.
  \item The Aether is not a simple weakly interacting Bose--Einstein condensate
    (which would have $K \propto n^2$), but a denser quantum liquid with a
    specific self-interaction potential.
  \item The healing length $\xi = \hbar/(m_A c \sqrt{2})$ sets the vortex core
    size and the minimum scale of the continuum description.
  \item Electron mass arises from vortex topology: $m_e c^2 = \frac{1}{2}\rhoa\,
    \kappa_e^2\,r_{\mathrm{loop}}[\ln(8r_{\mathrm{loop}}/\xi) - \frac{1}{2}]$, where $\kappa_e
    = h/(2m_A)$ is the half-quantum circulation.
  \item The Compton wavelength $\lambdabar_C = \hbar/(m_e c)$ is physically
    identified with the vortex core or loop scale, providing a physical basis
    for the QED ultraviolet cutoff.
  \item Length contraction arises microscopically from vortex geometry adjustment:
    $r_{\mathrm{loop}} \propto 1/\rhoa^2$ maintains constant rest energy as
    local density varies.
  \item The linearity constraint ($\delta\rho/\rhoa \ll 1$) favors $m_A \gg m_e$,
    corresponding to $\rhoa \gg 10^6$~kg/m$^3$ and an extremely stiff medium.
  \item Derrick's theorem proves that hedgehog skyrmions are unstable in the
    logarithmic condensate: the energy $E = AL + BL^3$ ($A, B > 0$) has no finite
    minimum, so the configuration collapses to a point.\ The electron must be a
    quantized vortex ring, stabilized by discrete circulation quantization
    $\kappa = nh/m_A$ rather than continuous energy balance.
\end{enumerate}




% ============================================================
\section{Experimental Support and Connections to Established Physics}
% ============================================================

All experimentally confirmed results of quantum mechanics---including the
Casimir effect, Lamb shift, anomalous magnetic moment, hydrogen spectroscopy,
Bell inequality violations, matter-wave interference, electron spin resonance,
and the photoelectric effect---are reproduced by this framework.\ The Aether
model adds no corrections to standard QM predictions at currently accessible
energies; deviations are predicted only at compressions approaching
$\delta\rho/\rhoa \sim 1$, far beyond experimental reach.


% ============================================================
\section{Open Questions and Testable Predictions}
% ============================================================

The following areas relevant to this paper remain underdeveloped and represent opportunities for further refinement or falsification:

\begin{enumerate}[leftmargin=2em]
  \item \textbf{Quantifying the DCE threshold for electron vortices.} The Dynamic Casimir Effect has
    well-defined mathematical conditions for when virtual particles become real.\ Mapping
    these conditions onto the electron spin vortex would yield the specific vortex velocity
    required for permanent binding, the expected ratio of transient to permanent events,
    and the energy cost of mass formation.\ This is the single most important formalization
    needed to make the theory quantitatively predictive.

  \item \textbf{Enhanced mass production in extreme environments.} The DCE mechanism predicts that
    environments with higher vortex velocities or Aether densities should produce more
    permanent binding events.\ This predicts enhanced particle production in neutron star
    interiors, near black hole event horizons, and in high-energy particle collisions---environments where effective boundary velocities approach or exceed the DCE threshold.

  \item \textbf{Schwinger, Unruh, and Hawking effect unification.} The Schwinger effect (pair
    production in strong electromagnetic fields), the Unruh effect (thermal radiation from
    acceleration), and Hawking radiation (thermal emission from black holes) all involve
    the conversion of vacuum fluctuations into observable particles or radiation.\ In the
    Aether model, all three arise from the same mechanism: disruption of the transient
    binding cycle's symmetry by velocity gradients or strong fields.\ The Schwinger
    threshold provides a field-strength condition, the Unruh temperature provides an
    acceleration condition, and the Hawking temperature provides a gravitational-gradient
    condition---all of which should be expressible in terms of Aether binding parameters.
    Demonstrating that all three effects yield consistent binding physics would strongly
    support the unified mechanism.

  \item \textbf{Quantitative unification of $G$ and $k$.} The theory predicts that the gravitational
    constant $G$ should be derivable from the Coulomb constant $k$, the speed of light $c$,
    and the asymmetry fraction $\epsilon$ of the transient bind--unbind cycle:
    $G = k \cdot \epsilon \cdot (\psi/\phi)^2 \times (\text{geometric factors})$.\ Calculating
    $\epsilon$ from the Heisenberg uncertainty principle and the DCE threshold would
    eliminate $G$ as an independent constant.

  \item \textbf{Quantitative derivation of the Aether flow profile near a nucleus.}\
    The resonant cavity model requires a specific flow gradient around the
    nucleus to produce the correct energy levels.\ Can this profile be derived
    self-consistently from the theory's existing description of electric fields as radial
    Aether inflow, combined with the gravitational inflow from the nuclear mass?\ The
    profile must yield both the Bohr radius ($a_0 = 0.529$~\AA) and the $1/n^2$ energy
    scaling.

  \item \textbf{Physical origin of electron spin quantization.}\ The resonant cavity
    model explains orbital quantization (angular phase matching) and energy quantization
    (radial confinement), but the quantization of electron spin---the fact that the spin
    vortex takes only two discrete orientations ($\pm\hbar/2$)---requires additional
    explanation.\ What property of the Aether vortex restricts it to exactly two states?


  \item \textbf{Derivation of anticommutation relations from Aether dynamics.}\
    The exclusion principle is mathematically encoded in the anticommutation
    relations $\{\hat{\psi}, \hat{\psi}^\dagger\} = 1$ for fermion field operators.\
    Can these be derived from the dynamical equations governing vortex interactions
    in the Aether, rather than imposed as an axiom?\ A successful derivation would
    demonstrate that quantum field theory's algebraic structure emerges from the
    medium's topology.

  \item \textbf{Boson--fermion distinction from vortex winding number.}\ The theory
    explains fermion exclusion through half-integer ($2\pi$) phase winding and boson
    non-exclusion through integer ($2\pi \times 2s$) winding.\ Can this be made
    quantitatively precise for composite particles?\ A deuteron (spin-1 boson) is
    composed of three spin-$\tfrac{1}{2}$ quarks---what is the vortex topology of the
    composite, and does it correctly predict bosonic statistics?


  \item \textbf{Photon polarization from condensate spin-orbit coupling.}\ The
    theory identifies photon polarization states with spin-wave excitations of
    the spinor condensate, with spatial transversality arising from spin-orbit
    coupling (Section~\ref{sec:qss}).\ Can the spin-orbit coupling Hamiltonian
    of the Aether condensate be specified, and does the resulting spin-wave
    dispersion relation reproduce Maxwell's equations---including the
    transversality condition $\mathbf{k} \cdot \mathbf{E} = 0$ and the
    absence of a longitudinal photon?

  \item \textbf{Pinning the Aether quantum mass $m_A$.}\ Three independent
    constraints (lattice QCD flux tube width, self-consistent vortex--bridge
    equations, and the glueball mass spectrum) converge on
    $m_A \approx 500$--$850$~MeV/$c^2$.\ Can precision measurements of vacuum
    fluctuations, gravitational wave stiffness, or next-order corrections to
    atomic spectra determine $m_A$ to better than $\pm 10\%$?

  \item \label{item:log-eos-micro}\textbf{Physical origin of the
    inverse-density interaction.}\ The $\Ka$-invariance constraint uniquely
    determines the Gross--Pitaevskii nonlinearity as $g(n) = \mu_0 - \Ka/n$
    (Eq.~\ref{eq:gp-nonlinearity})---an inverse-density interaction that is
    repulsive at low density and asymptotically constant at high density.\
    What microscopic physics between Aether quanta produces this specific
    functional form?\ Candidates include effective multi-body scattering
    at high density, quantum corrections to the condensate self-energy,
    or an intrinsic structural property of the Aether's internal degrees
    of freedom.\ Note that this is distinct from the logarithmic
    Gross--Pitaevskii potential $g(n) = b\ln(n/n_0)$
    \cite{rosen1969,zloshchastiev2011}, which yields a linear pressure
    $P = bn$ (density-independent speed, density-dependent bulk modulus)---the
    inverse of the Aether's requirements.

  \item \textbf{Breaking the $\rhoa$ degeneracy.}\ The vortex energy equation
    (\ref{eq:mass-from-topology}) connects $m_e$, $\rhoa$, and $m_A$ through
    the vortex geometry, but introduces $m_A$ as a new unknown.\ The system
    still has one degree of freedom.\ What additional constraint---from
    dispersion measurements, gravitational wave coupling, or vacuum
    impedance---can independently determine $\rhoa$ or $m_A$ and close the
    system?

  \item \textbf{Particle mass spectrum from vortex topology.}\ If the electron
    is a quantized vortex ring with circulation $\kappa_e = h/(2m_A)$ and loop
    radius $r_{\mathrm{loop}}$, can the full lepton mass spectrum (electron,
    muon, tau) be derived from the allowed quantized vortex configurations in
    the logarithmic condensate?\ Higher-energy vortex knots, linked rings, or
    multiply-wound configurations may correspond to heavier particles.\ Can the
    baryon masses (proton, neutron) similarly arise from three-vortex composite
    topologies?\ The Derrick instability of hedgehog skyrmions
    (Section~\ref{sec:derrick}) rules out one class of solitonic particles, but
    the full classification of stable vortex configurations in a spinor
    superfluid with logarithmic interactions remains open.

  \item \textbf{Formal derivation of Bell-violating correlations from the
    condensate.}\ The theory provides a physical mechanism for entanglement:
    correlated excitations of the quantum superfluid condensate inherit
    non-separable joint states from the condensate's global wavefunction,
    and the binding principle ensures that excitations cannot be described
    independently of the field they are bound to
    (Section~\ref{sec:entanglement-binding}).\ This mechanism is consistent
    with all observed entanglement phenomena and has a direct analogue in
    laboratory BEC pair production \cite{steinhauer2016}.\ What remains is
    a formal derivation: starting from the condensate's field operators and
    bosonic commutation relations, show that the resulting two-excitation
    correlation function reproduces exactly the $\cos^2(\theta_A - \theta_B)$
    dependence that violates Bell inequalities.\ This would elevate
    entanglement from a qualitative prediction to a quantitative one.

\end{enumerate}



% ============================================================
\section{Proposed Experiments}
% ============================================================

The following experiments target specific predictions relevant to this paper:

\begin{enumerate}[leftmargin=2em]
  \item \textbf{Quadratic Lorentz-violating photon dispersion.}\
    Special relativity predicts that Lorentz invariance is exact at all
    energies---photons of any frequency travel at exactly $c$.\ In the Aether
    model, Lorentz invariance is emergent and must break down at wavelengths
    approaching the relevant healing length.\ For photons (spin-wave
    excitations), the controlling scale is the \emph{spin} healing length
    $\xi_s$, which is much smaller than the density healing length $\xi$
    because the spin coupling vastly exceeds the density coupling
    (Section~\ref{sec:strong-constraints}).\ The predicted dispersion
    relation is $\omega^2 = k^2 c^2 (1 \pm k^2 \xi_s^2)$: a
    \emph{quadratic} energy dependence, not linear, because the isotropic
    medium preserves CPT symmetry.\ Current Fermi LAT constraints require
    $\xi_s \lesssim 10^{-27}$~m, consistent with the spin-wave channel
    being $\sim\!10^{22}$ times stiffer than the density channel.\
    The Cherenkov Telescope Array (CTA), with its improved sensitivity to
    time-of-arrival delays from high-redshift gamma-ray bursts, could push
    the quadratic LIV constraint by one to two orders of magnitude.\ A
    detection of quadratic (but not linear) dispersion would be a specific
    signature of a medium-based theory.

  \item \textbf{Vacuum birefringence from condensate anisotropy.}\
    Both photon polarizations are spin-wave excitations in the same
    channel, but they correspond to two transverse directions of the
    spinor oscillation.\ If the condensate ground state has any
    preferred orientation (e.g., residual nematic or ferromagnetic
    order), the two transverse directions become inequivalent and the
    polarization states propagate at marginally different speeds---vacuum
    birefringence.\ Standard electrodynamics (in the absence of external
    fields) predicts zero vacuum birefringence.\ The Aether model predicts
    $\Delta c / c \sim (E / E_{\mathrm{Planck}})^2$ from any small
    breaking of rotational symmetry within the spin-wave channel, which is
    undetectable at optical frequencies but could accumulate over cosmological
    distances for gamma-ray photons.\ Polarimetric observations of GRB
    afterglows and blazar jets already constrain vacuum birefringence at levels
    approaching $\Delta c / c < 10^{-38}$; future X-ray polarimetry missions
    (IXPE, eXTP) could tighten these bounds further.\ A null result would
    confirm that the condensate ground state is fully isotropic in the
    spin-wave sector.

\end{enumerate}

\noindent\textit{For all proposed experiments, see Paper~1: General Overview \cite{otto-overview}.}


% ============================================================
\section{Mathematical Summary}
% ============================================================

This section collects the key equations relevant to the topics covered in this paper.
\noindent\textit{For the complete mathematical summary, see Paper~1: General Overview \cite{otto-overview}.}

\subsection{Fundamental Aether Parameters}
% --------------------------------------------------

The theory posits four microscopic parameters from which all macroscopic
phenomena derive:
%
\begin{center}
\begin{tabular}{@{}llll@{}}
\hline
\textbf{Symbol} & \textbf{Name} & \textbf{Constrained range} & \textbf{Primary constraint} \\
\hline
$\rhoa$  & Ambient density     & $10^{11}$--$10^{12}$~kg/m$^3$         & Electron vortex energy \\
$\Ka$    & Bulk modulus         & $10^{28}$--$10^{29}$~J/m$^3$          & $\Ka = \rhoa\,c^2$ \\
$m_A$    & Aether quantum mass  & $500$--$850$~MeV/$c^2$                & Flux tube width / glueball mass \\
$\xi$    & Healing length       & $0.17$--$0.28$~fm                     & Bridge diameter (lattice QCD) \\
\hline
\end{tabular}
\end{center}

These four are not independent.\ Two defining relations connect them:
%
\begin{align}
  c &= \sqrt{\Ka / \rhoa}
    \label{eq:sum-speed-of-light} \\
  \xi &= \frac{\hbar}{m_A\,c\,\sqrt{2}}
    \label{eq:sum-healing-length}
\end{align}
%
leaving two free Aether parameters (e.g., $\rhoa$ and $m_A$) from which all
others follow.\ (A third parameter, $\epsilon$, enters through gravity; see below.)


% --------------------------------------------------
\subsection{Equation of State}
% --------------------------------------------------

The Aether obeys a logarithmic equation of state, the unique form that produces
a density-independent bulk modulus (Section~\ref{sec:qss}):
%
\begin{equation}
  P = \Ka \ln\!\Bigl(\frac{\rho}{\rho_0}\Bigr) + P_0
  \label{eq:sum-eos}
\end{equation}
%
Key consequence: the speed of sound $c_s = \sqrt{\Ka/\rho}$ varies with
density.\ At ambient density, $c_s = c$.\ Note: in gravitational fields,
Bernoulli's equation for free-fall inflow gives $\Delta\rho/\rhoa = 0$
exactly, so $c$ does not vary near gravitating masses; the density
dependence is relevant for material media and for the two-fluid
cosmological structure where the superfluid fraction varies with
temperature.\ This equation of state guarantees that $\Ka$ is constant
under compression, which is experimentally confirmed by:
%
\begin{itemize}[leftmargin=2em]
  \item Velocity time dilation measurements spanning $10^{-6}\,c$ to $0.9994\,c$
    (precision $\sim 10^{-8}$).
  \item The Fresnel drag coefficient $1 - 1/n^2$, exact across materials with
    density ratios up to $\sim$6:1.
\end{itemize}


% --------------------------------------------------
\subsection{Particle Properties from Vortex Topology}
% --------------------------------------------------

\paragraph{Circulation quantum.}
%
\begin{equation}
  \kappa_0 = \frac{h}{m_A}, \qquad
  \kappa_e = \frac{h}{2\,m_A} \quad (\text{half-quantum, electron})
  \label{eq:sum-circulation}
\end{equation}

\paragraph{Electron mass from vortex energy.}
%
\begin{equation}
  m_e\,c^2 = \tfrac{1}{2}\,\rhoa\,\kappa_e^2\,r_{\mathrm{loop}}\,
    \bigl[\ln\!\bigl(8r_{\mathrm{loop}}/\xi\bigr) - \tfrac{1}{2}\bigr]
  \label{eq:sum-electron-mass}
\end{equation}
%
with $r_{\mathrm{loop}} \approx \lambdabar_C = \hbar/(m_e c) \approx 386$~fm.

\paragraph{Vortex stability.}
Multiply quantized vortices ($n \geq 2$) are energetically unstable
($E \propto n^2$, while $n$ single vortices have $E \propto n$), providing
the Pauli exclusion principle.\ Each orbital mode admits at most two
opposite-circulation vortices.

\paragraph{Derrick's theorem.}
Hedgehog skyrmions in the logarithmic condensate have energy
$E = A L + B L^3$ with $A, B > 0$; no finite minimum exists.\ Particles must
be quantized vortex rings (Section~\ref{sec:derrick}).


% --------------------------------------------------
\subsection{Atomic Structure}
% --------------------------------------------------

\paragraph{Ground state energy.}
%
\begin{equation}
  E_1 = -\tfrac{1}{2}\,\alpha^2\,m_e\,c^2 = -13.6\;\mathrm{eV}
  \label{eq:sum-hydrogen}
\end{equation}

\paragraph{Quantization condition.}
Constructive self-interference in the resonant Aether cavity:
%
\begin{equation}
  n\,\lambda = 2\pi\,r
  \label{eq:sum-quantization}
\end{equation}

\paragraph{Fine structure constant.}
The electron's orbital velocity as a fraction of the Aether wave speed
(atomic Mach number):
%
\begin{equation}
  \alpha = \frac{v_{\mathrm{orbital}}}{c} \approx \frac{1}{137}
  \label{eq:sum-alpha}
\end{equation}


% --------------------------------------------------

\section*{Acknowledgments}
\addcontentsline{toc}{section}{Acknowledgments}

The author acknowledges the use of Claude (Anthropic) for \LaTeX{} formatting
and editorial assistance.

\bigskip


% ============================================================
% BIBLIOGRAPHY
% ============================================================

\begin{thebibliography}{99}

\bibitem{reinhardt2007}
Reinhardt, S., Saathoff, G., Buhr, H., et al.,
``Test of relativistic time dilation with fast optical atomic clocks at different velocities,''
\textit{Nature Physics}, \textbf{3}, 861--864 (2007).


\bibitem{duck1997}
Duck, I.\ and Sudarshan, E.C.G.,
\textit{Pauli and the Spin-Statistics Theorem}
(World Scientific, Singapore, 1997).


\bibitem{berry1997}
Berry, M.V.\ and Robbins, J.M.,
``Indistinguishability for quantum particles: spin, statistics and the geometric phase,''
\textit{Proc.\ R.\ Soc.\ Lond.\ A}, \textbf{453}, 1771--1790 (1997).


\bibitem{leinaas1977}
Leinaas, J.M.\ and Myrheim, J.,
``On the theory of identical particles,''
\textit{Il Nuovo Cimento B}, \textbf{37}, 1--23 (1977).


\bibitem{matthews1999}
Matthews, M.R.\ et al.,
``Vortices in a Bose--Einstein condensate,''
\textit{Phys.\ Rev.\ Lett.}, \textbf{83}, 2498--2501 (1999).


\bibitem{shin2004}
Shin, Y.\ et al.,
``Dynamical instability of a doubly quantized vortex in a Bose--Einstein condensate,''
\textit{Phys.\ Rev.\ Lett.}, \textbf{93}, 160406 (2004).


\bibitem{chandrasekhar1931}
Chandrasekhar, S.,
``The maximum mass of ideal white dwarfs,''
\textit{Astrophys.\ J.}, \textbf{74}, 81--82 (1931).


\bibitem{osheroff1972}
Osheroff, D.D., Richardson, R.C., and Lee, D.M.,
``Evidence for a new phase of solid He$^3$,''
\textit{Phys.\ Rev.\ Lett.}, \textbf{28}, 885--888 (1972).


\bibitem{autti2016}
Autti, S.\ et al.,
``Observation of half-quantum vortices in topological superfluid 3,''
\textit{Phys.\ Rev.\ Lett.}, \textbf{117}, 255301 (2016).



\bibitem{donnelly1991}
Donnelly, R.J.,
\textit{Quantized Vortices in Helium~II}
(Cambridge University Press, Cambridge, 1991).


\bibitem{vasileiou2013}
Vasileiou, V., Jacholkowska, A., Piron, F., et al.,
``Constraints on Lorentz invariance violation from Fermi--Large Area Telescope
observations of gamma-ray bursts,''
\textit{Phys.\ Rev.\ D}, \textbf{87}, 122001 (2013).


\bibitem{derrick1964}
Derrick, G.~H.,
``Comments on nonlinear wave equations as models for elementary particles,''
\textit{J.\ Math.\ Phys.}, \textbf{5}, 1252--1254 (1964).


\bibitem{skyrme1962}
Skyrme, T.~H.~R.,
``A unified field theory of mesons and baryons,''
\textit{Nucl.\ Phys.}, \textbf{31}, 556--569 (1962).


\bibitem{steinhauer2016}
Steinhauer, J.,
``Observation of quantum Hawking radiation and its entanglement in an analogue black hole,''
\textit{Nature Phys.}, \textbf{12}, 959--965 (2016).


\bibitem{volovik2003}
Volovik, G.~E.,
\textit{The Universe in a Helium Droplet}
(Oxford University Press, Oxford, 2003).


\bibitem{taylor1909}
Taylor, G.~I.,
``Interference fringes with feeble light,''
\textit{Proc.\ Camb.\ Phil.\ Soc.}, \textbf{15}, 114--115 (1909).


\bibitem{grangier1986}
Grangier, P., Roger, G., and Aspect, A.,
``Experimental evidence for a photon anticorrelation effect on a beam splitter: A new light on single-photon interferences,''
\textit{Europhys.\ Lett.}, \textbf{1}(4), 173--179 (1986).


\bibitem{fein2019}
Fein, Y.~Y., Geyer, P., Zwick, P., et~al.,
``Quantum superposition of molecules beyond 25\,kDa,''
\textit{Nat.\ Phys.}, \textbf{15}, 1242--1245 (2019).


\bibitem{wheeler1978}
Wheeler, J.~A.,
``The `past' and the `delayed-choice' double-slit experiment,''
in \textit{Mathematical Foundations of Quantum Theory}, ed.\ A.~R.\ Marlow
(Academic Press, New York, 1978), pp.\ 9--48.


\bibitem{jacques2007}
Jacques, V., Wu, E., Grosshans, F., Treussart, F., Grangier, P., Aspect, A., and Roch, J.-F.,
``Experimental realization of Wheeler's delayed-choice gedanken experiment,''
\textit{Science}, \textbf{315}(5814), 966--968 (2007).


\bibitem{elitzurvaidman1993}
Elitzur, A.~C.\ and Vaidman, L.,
``Quantum mechanical interaction-free measurements,''
\textit{Found.\ Phys.}, \textbf{23}(7), 987--997 (1993).


\bibitem{kwiat1995}
Kwiat, P., Weinfurter, H., Herzog, T., Zeilinger, A., and Kasevich, M.~A.,
``Interaction-free measurement,''
\textit{Phys.\ Rev.\ Lett.}, \textbf{74}(24), 4763--4766 (1995).


\bibitem{englert1996}
Englert, B.-G.,
``Fringe visibility and which-way information: An inequality,''
\textit{Phys.\ Rev.\ Lett.}, \textbf{77}(11), 2154--2157 (1996).


\bibitem{scullydruhl1982}
Scully, M.~O.\ and Dr\"uhl, K.,
``Quantum eraser: A proposed photon correlation experiment concerning observation and `delayed choice' in quantum mechanics,''
\textit{Phys.\ Rev.\ A}, \textbf{25}(4), 2208--2213 (1982).


\bibitem{kim2000}
Kim, Y.-H., Yu, R., Kulik, S.~P., Shih, Y., and Scully, M.~O.,
``Delayed `choice' quantum eraser,''
\textit{Phys.\ Rev.\ Lett.}, \textbf{84}(1), 1--5 (2000).


\bibitem{bell1964}
Bell, J.~S.,
``On the Einstein Podolsky Rosen paradox,''
\textit{Physics Physique Fizika}, \textbf{1}(3), 195--200 (1964).


\bibitem{stamper-kurn2013}
Stamper-Kurn, D.~M.\ and Ueda, M.,
``Spinor Bose gases: Symmetries, magnetism, and quantum dynamics,''
\textit{Rev.\ Mod.\ Phys.}, \textbf{85}, 1191--1244 (2013).


\bibitem{madelung1927}
Madelung, E.,
``Quantentheorie in hydrodynamischer Form,''
\textit{Z.\ Phys.}, \textbf{40}, 322--326 (1927).


\bibitem{wallstrom1994}
Wallstrom, T.~C.,
``Inequivalence between the Schr\"odinger equation and the Madelung hydrodynamic equations,''
\textit{Phys.\ Rev.\ A}, \textbf{49}(3), 1613--1617 (1994).


\bibitem{bogoliubov1947}
Bogoliubov, N.~N.,
``On the theory of superfluidity,''
\textit{J.\ Phys.\ (USSR)}, \textbf{11}, 23--32 (1947).


\bibitem{cardoso2013}
Cardoso, N., Cardoso, M., and Bicudo, P.,
``Inside the SU(3) quark-antiquark QCD flux tube: screening versus quantum widening,''
\textit{Phys.\ Rev.\ D}, \textbf{88}(5), 054504 (2013).


\bibitem{morningstar1999}
Morningstar, C.~J.\ and Peardon, M.~J.,
``The glueball spectrum from an anisotropic lattice study,''
\textit{Phys.\ Rev.\ D}, \textbf{60}(3), 034509 (1999).


\bibitem{cea2024}
Cea, P., Cosmai, L., and Ferretti, A.,
``The chromoelectric flux tube revisited,''
arXiv:2409.20168 [hep-lat] (2024).


\bibitem{cea2012}
Cea, P.\ and Cosmai, L.,
``Chromoelectric flux tubes and coherence length in QCD,''
\textit{Phys.\ Rev.\ D}, \textbf{86}(5), 054501 (2012).




\bibitem{otto-overview}
Otto, E.,
``Unifying Theory of Dynamic Aether Mechanics: General Overview,''
(2026).

\bibitem{otto-relativity}
Otto, E.,
``Unifying Theory of Dynamic Aether Mechanics: Special Relativity,''
(2026).

\bibitem{otto-cosmology}
Otto, E.,
``Unifying Theory of Dynamic Aether Mechanics: Cosmology,''
(2026).

\bibitem{otto-em}
Otto, E.,
``Unifying Theory of Dynamic Aether Mechanics: Magnetic and Electric Fields,''
(2026).

\bibitem{lorentz1904}
Lorentz, H.~A.,
``Electromagnetic phenomena in a system moving with any velocity smaller than that of light,''
\textit{Proceedings of the Royal Netherlands Academy of Arts and Sciences}, \textbf{6}, 809--831 (1904).

\bibitem{rosen1969}
Rosen, G.,
``Dilatation covariance and exact solutions in local relativistic field theories,''
\textit{J.\ Math.\ Phys.}, \textbf{10}, 1882--1887 (1969).

\bibitem{zloshchastiev2011}
Zloshchastiev, K.~G.,
``Volume element structure and roton-maxon-phonon excitations in superfluid helium beyond the Gross-Pitaevskii approximation,''
\textit{Eur.\ Phys.\ J.\ B}, \textbf{85}, 273 (2012).

\bibitem{bali2005}
Bali, G.~S., Neff, H., Duessel, T., Lippert, T., and Schilling, K.,
``Observation of string breaking in QCD,''
\textit{Phys.\ Rev.\ D}, \textbf{71}(11), 114513 (2005).

\bibitem{bulava2019}
Bulava, J., H\"orz, B., Knechtli, F., Koch, V., Moir, G., Morningstar, C., and Peardon, M.,
``String breaking by light and strange quarks in QCD,''
\textit{Phys.\ Lett.\ B}, \textbf{793}, 493--498 (2019).
\end{thebibliography}

\end{document}